\subsubsection{Oversampling}
\label{sec:semantics:grfrp:theorems:oversampling}

The execution model of \GRust presents a potential conflict between its two sub-languages. \GRFRP
services adopt a \emph{push-based} strategy, where computations are triggered by changes to imported
flows. In contrast, \GRSync components are compiled into \emph{pull-based} state machines that
conceptually expect to be evaluated at every logical \emph{step}. When a component is used within a
service, it is executed only when its inputs change. This mismatch could lead to incorrect behavior
if a component's state depends on the passage of time or internal feedback loops rather than on its
inputs.

To resolve this conflict, we introduced the reactive type system
(\Cref{sec:semantics:grsync:react_typing,sec:semantics:grfrp:react_typing}). This system enforces
that any component used within a service must be \emph{reactive}.
\Cref{thm:semantics:grsync:theorems:oversampling} states that a reactive component is one whose
internal state and outputs change only in response to changes in its \emph{discrete} inputs---those
whose values are determined by service imports and explicit timing constraints. This property
guarantees that an event-driven execution is sound.

The oversampling theorem (\Cref{thm:semantics:grfrp:theorems:oversampling}) formalizes this
guarantee at the service level. It states that evaluating a service at moments other than when its
imports change does not alter the values of its exported flows. Such additional evaluations are
referred to as \emph{oversampling}. This theorem provides the formal justification for shifting from
a traditional periodic execution model to a more efficient event-driven one without risking semantic
inconsistencies.

We model oversampling differently in our two semantic frameworks. In the denotational semantics,
oversampling corresponds to adding extra timestamps to the clock \tck, representing additional
evaluation instants. In the operational semantics, it is modeled by the arrival of an input message
from a flow not imported by the service. This technique simulates an external trigger for
computation that carries no new data, allowing us to observe the effect of an extra evaluation step
in isolation. The timestamp of this message propagates through the service \via the \timeop flow.

The proof of the oversampling theorem relies on \Cref{lem:semantics:grfrp:theorems:oversampling:op}
for flow operations. This lemma states that for any reactive operation, inserting an oversampling
step between two regular computation steps does not affect the final result. However, the 
operational semantics only makes flow \emph{updates} visible, where a $\bot$ value signifies no 
change. To formally reason about the outcome of a computation, we must first define the
\emph{current value} of a flow, which is the last non-$\bot$ value in its history of updates.

\newcommand{\CurrVal}[2]{\ensuremath{\smf{curr_{#1}}\left(#2\right)}}
We introduce the \smf{curr} function, which computes the current value of a flow from
the complete, unbounded list of its past values. This historical interpretation is distinct from the
memory-bounded \smf{into\_val} function of the operational semantics, which computes the next
state based only on the current state and a single update. The \smf{curr} function thus provides the
history-based view required for our proofs.

\begin{definition}[Current Value]
    \label{def:semantics:grfrp:theorems:oversampling:CurrVal}
    Consider a list $l$ of bottom values, representing the historical updates of a flow with type
    $sty$. The current value of the flow, denoted \CurrVal{sty}{l}, is defined recursively:
    \begin{align}
        \CurrVal{sty}{\emptylist}            & = v_d \label{eq:semantics:grfrp:theorems:oversampling:CurrVal:nil}                                      \\
        \CurrVal{sty}{l \lstconcat [v_\bot]} & = \intoVal{sty}{\CurrVal{sty}{l}}{v_\bot} \label{eq:semantics:grfrp:theorems:oversampling:CurrVal:cons}
    \end{align}
    with $v_d$ denoting a default value.

    For a signal of type $ty$ and a list of bottom values $l$, a bottom value $v_\bot$ represents an
    \emph{additional update} \iff $v_\bot \neq \bot$ implies $\CurrVal{ty}{l} \neq v_\bot$. We say
    that $l$ is \emph{a list of additional updates} if, for any element $v_\bot$ in $l$, $v_\bot$ is
    an additional update with respect to the sub-list of elements in $l$ that precede $v_\bot$.
\end{definition}

\begin{remark}[Propagation of Additional Updates]
    The property of being an \emph{additional update} is preserved across reactive flow
    computations. We observe that if all dependencies to a flow operation are additional updates,
    its resulting values will also be additional updates.
    %
    This observation extends to the service level: if the imported flow values are additional
    updates, then all intermediate flows computed within the service will also consist of additional
    updates.
    %
    While a formal proof is not provided in this manuscript, we conjecture that these properties can
    be established by structural induction over the definitions of flow operations and service
    statements.
\end{remark}

\begin{lemma}[Current value equalities]
    \label{lem:semantics:grfrp:theorems:oversampling:CurrVal:last}
    For all flow of type $sty$, all list bottom of values $l$, and all bottom values $v_{1\bot}$ and
    $v_{2\bot}$ representing additional updates of $l$, we have:
    \begin{equation*}
        \CurrVal{sty}{l \lstconcat [v_{1\bot}]} = \CurrVal{sty}{l \lstconcat [v_{2\bot}]}
        \Leftrightarrow
        v_{1\bot} = v_{2\bot}.
    \end{equation*}
\end{lemma}
\begin{skippableproof}
    We establish this equivalence by proving implication in both directions.

    \noindent\textit{Case $\Leftarrow$.}
    %
    Let a flow of type $sty$, $l$ be a list bottom of values, and $v_{1\bot}$ and $v_{2\bot}$ be two
    bottom values representing additional updates of $l$ such that $v_{1\bot} = v_{2\bot}$. The
    function \CurrVal{sty}{\_} is deterministic, so that we have
    $\CurrVal{sty}{l \lstconcat [v_{1\bot}]} = \CurrVal{sty}{l \lstconcat [v_{2\bot}]}$.

    \noindent\textit{Case $\Rightarrow$.}
    %
    Let a flow of type $sty$, $l$ be a list bottom of values, and $v_{1\bot}$ and $v_{2\bot}$ be two
    bottom values representing additional updates of $l$ such that:
    \begin{equation*}
        \CurrVal{sty}{l \lstconcat [v_{1\bot}]} = \CurrVal{sty}{l \lstconcat [v_{2\bot}]}.
    \end{equation*}

    \begin{itemize}
        \item \textit{If there exist a type $ty$ such that $sty = \mayB{ty}$
                  (\ie it is an event type):}

              \Cref{eq:semantics:grfrp:theorems:oversampling:CurrVal:cons,%
                  rule:semantics:grfrp:opsem:aux:intoVal:event} ensure that for
              $i \in \{1,2\}$ we have:
              \begin{align*}
                  \CurrVal{\mayB{ty}}{l \lstconcat [v_{i\bot}]} & = \intoVal{\mayB{ty}}{\CurrVal{\mayB{ty}}{l}}{v_{i\bot}}                                     \\
                                                                & = \ttf{match} \, v_{i\bot} \, \{ \, v \rightarrow \some{v}, \, \bot \rightarrow \none \, \}.
              \end{align*}

              Then, if \CurrVal{\mayB{ty}}{l \lstconcat [v_{i\bot}]} are equal to \none, then both
              $v_{1\bot}$ and $v_{2\bot}$ are equal to $\bot$.

              Conversely, if there exist a value $v$ such that
              \CurrVal{\mayB{ty}}{l \lstconcat [v_{i\bot}]} are equal to \some{v}, then both
              $v_{1\bot}$ and $v_{2\bot}$ are equal to $v$.

        \item \textit{If there exist a type $ty$ such that $sty = ty$ (\ie it is a signal type):}

              \Cref{eq:semantics:grfrp:theorems:oversampling:CurrVal:cons,%
                  rule:semantics:grfrp:opsem:aux:intoVal:signal} ensure that for
              $i \in \{1,2\}$ we have:
              \begin{align*}
                  \CurrVal{ty}{l \lstconcat [v_{i\bot}]} & = \intoVal{ty}{\CurrVal{ty}{l}}{v_{i\bot}}                                                      \\
                                                         & = \ttf{match} \, v_{i\bot} \, \{ \, v \rightarrow v, \, \bot \rightarrow \CurrVal{ty}{l} \, \}.
              \end{align*}

              Then, if \CurrVal{ty}{l \lstconcat [v_{i\bot}]} are equal to \CurrVal{ty}{l}, because
              $\CurrVal{ty}{l} \neq v_{i\bot}$, then both $v_{1\bot}$ and $v_{2\bot}$ are equal to
              $\bot$.

              Conversely, if there exist a value $v$ such that \CurrVal{ty}{l \lstconcat [v_{i\bot}]}
              are equal to $v$, because $\CurrVal{ty}{l} \neq v_{i\bot}$, then both $v_{1\bot}$ and
              $v_{2\bot}$ are equal to $v$.
    \end{itemize}

    In all cases, we have $v_{1\bot} = v_{2\bot}$.
\end{skippableproof}

With the \smf{curr} function, we can precisely formulate the guarantee provided by the oversampling
\Cref{lem:semantics:grfrp:theorems:oversampling:op} for flow operations. Consider a simulation of
$N$ steps. If we perform an oversampling evaluation just after the $n^\text{th}$ step
(with $n < N$), the system enters an intermediate state. The lemma guarantees that this intermediate
step has no lasting effect. Subsequently, if the $(n+1)^\text{th}$ step is performed from this state
using the original inputs, the resulting current value of any flow is identical to what it would
have been in the non-oversampled simulation.

\begin{lemma}[Oversampling of \GRFRP reactive flow operations]
    \label{lem:semantics:grfrp:theorems:oversampling:op}
    Let $\sigma$ be a flow operation in a program \G and typing context \tyCtx, let $sty^+$ be
    types, let $\rType^+$ be reactive types, and let $\Theta$ be a reactive typing context such
    that:
    \begin{equation*}
        \flowTyped{\sigma}{sty^+} \qquad \opTypedReact{\Theta}{\sigma}{\rType^+}.
    \end{equation*}
    For
    all $n \ge 1$;
    all states $s_0$, \dots, $s_{n+1}$, $s'_{n+1}$, and $s'_{n+2}$ corresponding to $\sigma$;
    all update context $\gamma_0$, \dots, $\gamma_n$, $\gamma'_n$, and $\gamma'_{n+1}$;
    all timestamps $t_0$, \dots, $t_n$, $t'_n$, and $t'_{n+1}$;
    all bottom values $v_{0\bot}^+$, \dots, $v_{n\bot}^+$, $v_{n\bot}^{\prime+}$, and $v_{n+1\bot}^{\prime+}$; and
    all list of timer resets $[(i_{T0}, t_{T0})^+]$, \dots, $[(i_{Tn}, t_{Tn})^+]$, $[(i'_{Tn}, t'_{Tn})^+]$, and $[(i'_{Tn+1}, t'_{Tn+1})^+]$
    such that:
    \begin{equation*}
        \begin{array}{c}
            t'_n \le t'_{n+1} = t_n \qquad s_0 = \initstate{\sigma}                                                                \\
            \forall k \in \{0, \dots, n\}, \quad \sigmaOpSem{\gamma_k}{s_k}{\sigma}{t_k}{s_{k+1}}{v_{k\bot}^+}{(i_{Tk}, t_{Tk})^+} \\
            \sigmaOpSem{\gamma'_n}{s_n}{\sigma}{t'_n}{s'_{n+1}}{v_{n\bot}^{\prime+}}{(i'_{Tn}, t'_{Tn})^+}                         \\
            \sigmaOpSem{\gamma'_{n+1}}{s'_{n+1}}{\sigma}{t'_{n+1}}{s'_{n+2}}{v_{n+1\bot}^{\prime+}}{(i'_{Tn+1}, t'_{Tn+1})^+},
        \end{array}
    \end{equation*}
    we have:
    \begin{multline*}
        \left[\begin{array}{l}
                \forall x \in \depSet{\sigma}, \; \left(\Theta[x] = \dType \Rightarrow \gamma'_n[x] = \bot\right) \\
                \quad \wedge \left(\CurrVal{\tyCtx[x]}{[\gamma_0[x], \dots, \gamma_n[x]]} = \CurrVal{\tyCtx[x]}{[\gamma_0[x], \dots, \gamma_{n-1}[x], \gamma'_n[x], \gamma'_{n+1}[x]]}\right)
            \end{array}\right]
        \quad \Rightarrow \\
        \begin{array}{ll}
                   & s'_{n+2} = s_{n+1} \wedge [(i'_{Tn}, t'_{Tn})^+] = \emptylist \wedge [(i'_{Tn+1}, t'_{Tn+1})^+] = [(i_{Tn}, t_{Tn})^+]                  \\
            \wedge & \left(\rType = \dType \Rightarrow v'_{n\bot} = \bot\right)^+                                                                            \\
            \wedge & \left(\CurrVal{sty}{[v_{0\bot}, \dots, v_{n\bot}]} = \CurrVal{sty}{[v_{0\bot}, \dots, v_{n-1\bot}, v'_{n\bot}, v'_{n+1\bot}]}\right)^+.
        \end{array}
    \end{multline*}

    This lemma formalizes the behavior of a flow operation, $\sigma$, during an oversampling step.
    It states that if, after $n$ simulation steps, all discrete dependencies of $\sigma$
    (of type \dType) are oversampled at step $n+1$, then this intermediate step has no lasting
    effect. Specifically, a computation performed from this new state yields the same result as
    one performed from the state at step $n$. This implies that the internal state of the flow
    operation does not change ($s'_{n+2} = s_{n+1}$), timer resets are preserved
    ($[(i'_{Tn}, t'_{Tn})^+] = \emptylist$ and $[(i'_{Tn+1}, t'_{Tn+1})^+] = [(i_{Tn}, t_{Tn})^+]$),
    discrete outputs are oversampled ($v'_{n\bot} = \bot$), and the current value of the outputs
    remains identical ($\CurrVal{sty}{[v_{0\bot}, \dots, v_{n\bot}]} =
        \CurrVal{sty}{[v_{0\bot}, \dots, v_{n-1\bot}, v'_{n\bot}, v'_{n+1\bot}]}$).
\end{lemma}
\begin{skippableproof}
    Let $\sigma$ be a flow operation in a program \G and typing context \tyCtx, let $sty^+$ be
    types, let $\rType^+$ be reactive types, let $\Theta$ be a reactive typing context, let $n\ge 1$,
    let $s_0$, \dots, $s_{n+1}$, $s'_{n+1}$, and $s'_{n+2}$ be states corresponding to $\sigma$, let
    $\gamma_0$, \dots, $\gamma_n$, $\gamma'_n$, and $\gamma'_{n+1}$ be update context, let $t_0$,
    \dots, $t_n$, $t'_n$, and $t'_{n+1}$ be timestamps, let $v_{0\bot}^+$, \dots, $v_{n\bot}^+$,
    $v_{n\bot}^{\prime+}$, and $v_{n+1\bot}^{\prime+}$ be bottom values, and let
    $[(i_{T0}, t_{T0})^+]$, \dots, $[(i_{Tn}, t_{Tn})^+]$, $[(i'_{Tn}, t'_{Tn})^+]$, and
    $[(i'_{Tn+1}, t'_{Tn+1})^+]$ be list of timer resets such that:
    \begin{equation*}
        \begin{array}{c}
            t'_n \le t'_{n+1} = t_n \qquad s_0 = \initstate{\sigma}                                                                \\
            \forall k \in \{0, \dots, n\}, \quad \sigmaOpSem{\gamma_k}{s_k}{\sigma}{t_k}{s_{k+1}}{v_{k\bot}^+}{(i_{Tk}, t_{Tk})^+} \\
            \sigmaOpSem{\gamma'_n}{s_n}{\sigma}{t'_n}{s'_{n+1}}{v_{n\bot}^{\prime+}}{(i'_{Tn}, t'_{Tn})^+}                         \\
            \sigmaOpSem{\gamma'_{n+1}}{s'_{n+1}}{\sigma}{t'_{n+1}}{s'_{n+2}}{v_{n+1\bot}^{\prime+}}{(i'_{Tn+1}, t'_{Tn+1})^+},
        \end{array}
    \end{equation*}
    and,
    \begin{multline*}
        \forall x \in \depSet{\sigma}, \; \left(\Theta[x] = \dType \Rightarrow \gamma'_n[x] = \bot\right) \\
        \wedge \left(\CurrVal{\tyCtx[x]}{[\gamma_0[x], \dots, \gamma_n[x]]} = \CurrVal{\tyCtx[x]}{[\gamma_0[x], \dots, \gamma_{n-1}[x], \gamma'_n[x], \gamma'_{n+1}[x]]}\right)
    \end{multline*}

    We proceed by structural induction on $\sigma$.

    \paragraph{Identifiers}
    Suppose $\sigma$ is $x$, the identifier to the flow $x$.

    The reactive typing rule \nameref{rule:semantics:grfrp:react_typing:op:ident} ensures that
    $\rType$ equals $\Theta[x]$. The identifier $x$ is in \depSet{x}, ensured by
    \Cref{eq:semantics:grfrp:theorems:determinism:deps:ident}, which implies that the current value
    of $x$ in the context without oversampling (\CurrVal{\tyCtx[x]}{[\gamma_0[x], \dots, \gamma_n[x]]}) equals
    its current value in the context with oversampling (\CurrVal{\tyCtx[x]}{[\gamma_0[x], \dots, \gamma_{n-1}[x],
    \gamma'_n[x], \gamma'_{n+1}[x]]}) and if $\Theta[x]$ equals $\dType$ then $\gamma'_n[x]$ is
    $\bot$.

    The typing rule \nameref{rule:grust:typing:grfrp:flow_op:ident} ensures that $sty = \tyCtx[x]$.

    The operational semantics rule \nameref{rule:semantics:grfrp:opsem:step:op:ident} ensures that
    $s'_{n+2} = ()$ and $s_{n+1} = ()$, that $v_{n\bot} = \gamma_n[x]$, $v'_{n\bot} = \gamma'_n[x]$,
    and $v'_{n+1\bot} = \gamma'_{n+1}[x]$, and that $[(i_{Tn}, t_{Tn})^+] = [(i'_{Tn}, t'_{Tn})^+] =
        [(i'_{Tn+1}, t'_{Tn+1})^+] = \emptylist$.

    It follows that $s'_{n+2} = s_{n+1}$, that the current value of $x$
    (\CurrVal{sty}{[v_{0\bot}, \dots, v_{n\bot}]}) remains the same with oversampling (\CurrVal{sty}{[v_{0\bot},
    \dots, v_{n-1\bot}, v'_{n\bot}, v'_{n+1\bot}]}), that if $\rType = \dType$ then $v'_{n\bot}$ is
    $\bot$, that $[(i'_{Tn}, t'_{Tn})^+]$ is empty, and that $[(i'_{Tn+1}, t'_{Tn+1})^+]$ equals
    $[(i_{Tn}, t_{Tn})^+]$.

    \paragraph{Time}
    Suppose $\sigma$ is \timeop, the time operation.

    The operational semantics rule \nameref{rule:semantics:grfrp:opsem:step:op:time} ensures that:
    \begin{equation*}
        \begin{array}{cc}
            s_0 = ()      & \sigmaOpSem{\gamma_0}{()}{\timeop}{t_0}{()}{t_0}{\,}                 \\
            \dots         & \dots                                                                \\
            s_{n+1} = ()  & \sigmaOpSem{\gamma_n}{()}{\timeop}{t_n}{()}{t_n}{\,}                 \\
            s'_{n+1} = () & \sigmaOpSem{\gamma'_n}{()}{\timeop}{t'_n}{()}{t'_n}{\,}              \\
            s'_{n+2} = () & \sigmaOpSem{\gamma'_{n+1}}{()}{\timeop}{t'_{n+1}}{()}{t'_{n+1}}{\,}.
        \end{array}
    \end{equation*}

    The reactive typing rule \nameref{rule:semantics:grfrp:react_typing:op:time} ensures that
    $\rType = \cType$.

    \begin{align*}
        [(i'_{Tn}, t'_{Tn})^+]                                                   & = \emptylist                                           \\
        s'_{n+2}                                                                 & = () = s_{n+1}                                         \\
        \CurrVal{sty}{[v_{0\bot}, \dots, v_{n-1\bot}, v'_{n\bot}, v'_{n+1\bot}]} & = \CurrVal{sty}{[t_0, \dots, t_{n-1}, t'_n, t'_{n+1}]} \\
                                                                                 & = t'_{n+1} = t_n = \CurrVal{sty}{[t_0, \dots, t_n]}    \\
                                                                                 & = \CurrVal{sty}{[v_{0\bot}, \dots, v_{n\bot}]}         \\
        [(i'_{Tn+1}, t'_{Tn+1})^+]                                               & = \emptylist = [(i_{Tn}, t_{Tn})^+],
    \end{align*}
    which concludes the proof for the \timeop operator.

    \paragraph{Sample}
    Suppose $\sigma$ is \sample{\sigma}{T}, the sample of the event $\sigma$ on the period $T$.

    The operational semantics rules \nameref{rule:semantics:grfrp:opsem:step:op:sample:propag},
    \nameref{rule:semantics:grfrp:opsem:step:op:sample:timer} and
    \nameref{rule:semantics:grfrp:opsem:step:op:sample:value} ensure there exist states $s_{\sigma 0}$,
    \dots, $s_{\sigma n+1}$, $s'_{\sigma n+1}$ and $s'_{\sigma n+2}$ corresponding to the input $\sigma$, bottom
    values $v_{\sigma 0\bot}$, \dots, $v_{\sigma n\bot}$, $v'_{\sigma n\bot}$ and $v'_{\sigma n+1\bot}$, timer resets
    $[(i_{\sigma0}, t_{\sigma0})^+]$, \dots, $[(i_{\sigma n}, t_{\sigma n})^+]$, $[(i'_{\sigma n},
        t'_{\sigma n})^+]$ and $[(i'_{\sigma n+1}, t'_{\sigma n+1})^+]$, and bottom values
    $(v_{m0\bot})$, \dots, $(v_{mn+1\bot})$, $(v'_{mn+1\bot})$, $(v'_{mn+2\bot})$,
    such that:
    \begin{equation*}
        \begin{array}{cc}
            s_0 = (v_{m0\bot}, s_{\sigma 0})            & \sigmaOpSem{\gamma_0}{s_{\sigma 0}}{\sigma}{t_0}{s_{\sigma 1}}{v_{\sigma 0\bot}}{(i_{\sigma0}, t_{\sigma0})^+}                             \\
            \dots                                       & \dots                                                                                                                                      \\
            s_{n+1} = (v_{mn+1\bot}, s_{\sigma n+1})    & \sigmaOpSem{\gamma_n}{s_{\sigma n}}{\sigma}{t_n}{s_{\sigma n+1}}{v_{\sigma n\bot}}{(i_{\sigma n}, t_{\sigma n})^+}                         \\
            s'_{n+1} = (v'_{mn+1\bot}, s'_{\sigma n+1}) & \sigmaOpSem{\gamma'_n}{s_{\sigma n}}{\sigma}{t'_n}{s'_{\sigma n+1}}{v'_{\sigma n\bot}}{(i'_{\sigma n}, t'_{\sigma n})^+}                   \\
            s'_{n+2} = (v'_{mn+2\bot}, s'_{\sigma n+2}) & \sigmaOpSem{\gamma'_{n+1}}{s'_{\sigma n+1}}{\sigma}{t'_{n+1}}{s'_{\sigma n+2}}{v'_{\sigma n+1\bot}}{(i'_{\sigma n+1}, t'_{\sigma n+1})^+}.
        \end{array}
    \end{equation*}

    The typing rule \nameref{rule:grust:typing:grfrp:flow_op:sample} ensures there exists a types
    $ty_\sigma$ such that:
    \begin{align*}
        \flowTyped{\sigma}{\mayB{ty_\sigma}}.
    \end{align*}

    The reactive typing rule \nameref{rule:semantics:grfrp:react_typing:op:sample} ensures that
    $\rType = \dType$ and:
    \begin{equation*}
        \opTypedReact{\Theta}{\sigma}{\dType}.
    \end{equation*}

    \Cref{eq:semantics:grfrp:theorems:determinism:deps:sample} ensures that for all $x$ in
    \depSet{\sigma}, we have $x$ in \depSet{\sample{\sigma}{T}}, therefore the current value of $x$ in
    the context without oversampling (encoded by the value \CurrVal{\tyCtx[x]}{[\gamma_0[x], \dots, \gamma_n[x]]})
    equals its current value in the context with oversampling (encoded by
    \CurrVal{\tyCtx[x]}{[\gamma_0[x], \dots, \gamma_{n-1}[x], \gamma'_n[x], \gamma'_{n+1}[x]]}), and
    if $\Theta[x]$ equals $\dType$ then $\gamma'_n[x]$ is $\bot$.

    It is then possible to apply the structural induction hypothesis on $\sigma$, which implies that
    $s'_{\sigma n+2} = s_{\sigma n+2}$, $[(i'_{\sigma n}, t'_{\sigma n})^+]$ is empty, that the
    timer resets $[(i'_{\sigma n+1}, t'_{\sigma n+1})^+]$ equals $[(i_{\sigma n}, t_{\sigma n})^+]$,
    that the current value of the value of $\sigma$ without oversampling (\CurrVal{\mayB{ty_\sigma}}{[v_{\sigma 0\bot},
                \dots, v_{\sigma n\bot}]}) equals the one with oversampling (\CurrVal{\mayB{ty_\sigma}}{[v_{\sigma 0\bot}, \dots,
    v_{\sigma n-1\bot}, v'_{\sigma n\bot}, v'_{\sigma n+1\bot}]}), and that $v'_{\sigma n\bot}$ is $\bot$.

    The presence of a timer identifier ($i_T$) in an update context ($\gamma$) can only come from
    the propagation of a timer message (\timerMess{i_T}), as it is only shown by the rule
    \nameref{rule:semantics:grfrp:opsem:step:service:timer}.
    Therefore, timer identifiers are not present in update contexts during the oversampling at time
    $t'_n$. The only possible step transition is the rule
    \nameref{rule:semantics:grfrp:opsem:step:op:sample:propag} ($v'_{\sigma n\bot} = \bot$ disable
    the rule \nameref{rule:semantics:grfrp:opsem:step:op:sample:value}).

    The rule \nameref{rule:semantics:grfrp:opsem:step:op:sample:propag} ensures that:\footnote{The notation \myeq{IH} denotes the application of the induction hypothesis to establish the equality.}
    \begin{align*}
        s'_{n+1}               & = (v_{mn\bot}, s'_{\sigma n+1})                            \\
        v'_{n\bot}             & = \bot                                                     \\
        [(i'_{Tn}, t'_{Tn})^+] & = [(i'_{\sigma n}, t'_{\sigma n})^+] \myeq{IH} \emptylist.
    \end{align*}

    We can now prove that, in any transition case, the step after oversampling is equivalent to the
    step without oversampling.
    \begin{itemize}
        \item Suppose there exist a value $v$ such that $v'_{\sigma n+1\bot} = v$.
              The induction hypothesis ensures:
              \begin{align*}
                  \CurrVal{\mayB{ty_\sigma}}{[v_{\sigma 0\bot}, \dots, v_{\sigma n\bot}]} & = \CurrVal{\mayB{ty_\sigma}}{[v_{\sigma 0\bot}, \dots, v_{\sigma n-1\bot}, v'_{\sigma n\bot}, v'_{\sigma n+1\bot}]} \\
                                                                                          & = \CurrVal{\mayB{ty_\sigma}}{[v_{\sigma 0\bot}, \dots, v_{\sigma n-1\bot}, \bot, v'_{\sigma n+1\bot}]}              \\
                                                                                          & = \CurrVal{\mayB{ty_\sigma}}{[v_{\sigma 0\bot}, \dots, v_{\sigma n-1\bot}, v'_{\sigma n+1\bot}]}                    \\
                                                                                          & \qquad \text{\cf \Cref{rule:semantics:grfrp:opsem:aux:intoVal:event}}                                               \\
                                                                                          & \Rightarrow \quad v_{\sigma n\bot} = v'_{\sigma n+1\bot} = v                                                        \\
                                                                                          & \qquad \text{\cf \Cref{lem:semantics:grfrp:theorems:oversampling:CurrVal:last}}.
              \end{align*}

              The rule \nameref{rule:semantics:grfrp:opsem:step:op:sample:value} can then be applied
              on both steps (the one at $t_n$ and the one at $t'_{n+1}$) and ensures that:
              \begin{align*}
                  s'_{n+2}                   & = (v, s'_{\sigma n+2}) \myeq{IH} (v, s_{\sigma n+1})                                \\
                  s_{n+1}                    & = (v, s_{\sigma n+1})                                                               \\
                  v'_{n+1\bot}               & = \bot                                                                              \\
                  v_{n\bot}                  & = \bot                                                                              \\
                  [(i'_{Tn+1}, t'_{Tn+1})^+] & = [(i'_{\sigma n+1}, t'_{\sigma n+1})^+] \myeq{IH} [(i_{\sigma n}, t_{\sigma n})^+] \\
                  [(i_{Tn}, t_{Tn})^+]       & = [(i_{\sigma n}, t_{\sigma n})^+],
              \end{align*}
              so that the current values are:
              \begin{align*}
                  \CurrVal{ty_\sigma}{[v_{0\bot}, \dots, v_{n-1\bot}, v_{n\bot}]}                & = \CurrVal{ty_\sigma}{[v_{0\bot}, \dots, v_{n-1\bot}, \bot]}            \\
                                                                                                 & = \CurrVal{ty_\sigma}{[v_{0\bot}, \dots, v_{n-1\bot}]}                  \\
                                                                                                 & \qquad \text{\cf \Cref{rule:semantics:grfrp:opsem:aux:intoVal:signal}}  \\
                  \CurrVal{ty_\sigma}{[v_{0\bot}, \dots, v_{n-1\bot}, v'_{n\bot}, v'_{n+1\bot}]} & = \CurrVal{ty_\sigma}{[v_{0\bot}, \dots, v_{n-1\bot}, \bot, \bot]}      \\
                                                                                                 & = \CurrVal{ty_\sigma}{[v_{0\bot}, \dots, v_{n-1\bot}]}                  \\
                                                                                                 & \qquad \text{\cf \Cref{rule:semantics:grfrp:opsem:aux:intoVal:signal}},
              \end{align*}
              which implies:
              \begin{align*}
                  s'_{n+2}                                                                       & = s_{n+1}                                            \\
                  \CurrVal{ty_\sigma}{[v_{0\bot}, \dots, v_{n-1\bot}, v'_{n\bot}, v'_{n+1\bot}]} & = \CurrVal{ty_\sigma}{[v_{0\bot}, \dots, v_{n\bot}]} \\
                  [(i'_{Tn+1}, t'_{Tn+1})^+]                                                     & = [(i_{Tn}, t_{Tn})^+].
              \end{align*}

        \item Suppose $v'_{\sigma n+1\bot} = \bot$ and $\gamma'_{n+1}[i_T] \neq \bot$ (meaning there
              is a timer at time $t'_{n+1}$). Because $t'_{n+1} = t_n$ and that \emph{no timers were
                  reset during the oversampling} at time $t'_n$, the occurrence of the timer at
              $t'_{n+1}$ implies the occurrence of the same timer at time $t_n$
              ($\gamma_n[i_T] \neq \bot$).

              The rule \nameref{rule:semantics:grfrp:opsem:step:op:sample:timer} can then be applied
              on both steps and ensures that:
              \begin{align*}
                  s'_{n+2}                   & = (\bot, s'_{\sigma n+2}) \myeq{IH} (\bot, s_{\sigma n+1})                                                               \\
                  s_{n+1}                    & = (\bot, s_{\sigma n+1})                                                                                                 \\
                  v'_{n+1\bot}               & = v'_{mn+1\bot} = v_{mn\bot}                                                                                             \\
                  v_{n\bot}                  & = v_{mn\bot}                                                                                                             \\
                  [(i'_{Tn+1}, t'_{Tn+1})^+] & = [(i_T, t'_{n+1} + T), (i'_{\sigma n+1}, t'_{\sigma n+1})^+] \myeq{IH} [(i_T, t_n + T), (i_{\sigma n}, t_{\sigma n})^+] \\
                  [(i_{Tn}, t_{Tn})^+]       & = [(i_T, t_n + T), (i_{\sigma n}, t_{\sigma n})^+],
              \end{align*}
              so that the current values are:
              \begin{align*}
                  \CurrVal{ty_\sigma}{[v_{0\bot}, \dots, v_{n-1\bot}, v_{n\bot}]}                & = \CurrVal{ty_\sigma}{[v_{0\bot}, \dots, v_{n-1\bot}, v_{mn\bot}]}       \\
                  \CurrVal{ty_\sigma}{[v_{0\bot}, \dots, v_{n-1\bot}, v'_{n\bot}, v'_{n+1\bot}]} & = \CurrVal{ty_\sigma}{[v_{0\bot}, \dots, v_{n-1\bot}, \bot, v_{mn\bot}]} \\
                                                                                                 & = \CurrVal{ty_\sigma}{[v_{0\bot}, \dots, v_{n-1\bot}, v_{mn\bot}]}       \\
                                                                                                 & \qquad \text{\cf \Cref{rule:semantics:grfrp:opsem:aux:intoVal:signal}},
              \end{align*}
              which implies:
              \begin{align*}
                  s'_{n+2}                                                                       & = s_{n+1}                                            \\
                  \CurrVal{ty_\sigma}{[v_{0\bot}, \dots, v_{n-1\bot}, v'_{n\bot}, v'_{n+1\bot}]} & = \CurrVal{ty_\sigma}{[v_{0\bot}, \dots, v_{n\bot}]} \\
                  [(i'_{Tn+1}, t'_{Tn+1})^+]                                                     & = [(i_{Tn}, t_{Tn})^+].
              \end{align*}

        \item Suppose $v'_{\sigma n+1\bot} = \bot$ and $\gamma'_{n+1}[i_T] = \bot$. It is the same
              argument as in the previous point: because $t'_{n+1} = t_n$ and that no timers were
              reset during the oversampling at time $t'_n$, the absence of the timer at $t'_{n+1}$
              implies the absence of the same timer at time $t_n$ ($\gamma_n[i_T] = \bot$).

              The rule \nameref{rule:semantics:grfrp:opsem:step:op:sample:propag} can then be
              applied on both steps and ensures that:
              \begin{align*}
                  s'_{n+2}                   & = (v'_{mn+1\bot}, s'_{\sigma n+2}) = (v_{mn\bot}, s'_{\sigma n+2}) \myeq{IH} (v_{mn\bot}, s_{\sigma n+1}) \\
                  s_{n+1}                    & = (v_{mn\bot}, s_{\sigma n+1})                                                                            \\
                  v'_{n+1\bot}               & = \bot                                                                                                    \\
                  v_{n\bot}                  & = \bot                                                                                                    \\
                  [(i'_{Tn+1}, t'_{Tn+1})^+] & = [(i'_{\sigma n+1}, t'_{\sigma n+1})^+] \myeq{IH} [(i_{\sigma n}, t_{\sigma n})^+]                       \\
                  [(i_{Tn}, t_{Tn})^+]       & = [(i_{\sigma n}, t_{\sigma n})^+],
              \end{align*}
              so that the current values are:
              \begin{align*}
                  \CurrVal{ty_\sigma}{[v_{0\bot}, \dots, v_{n-1\bot}, v_{n\bot}]}                & = \CurrVal{ty_\sigma}{[v_{0\bot}, \dots, v_{n-1\bot}, \bot]}            \\
                                                                                                 & = \CurrVal{ty_\sigma}{[v_{0\bot}, \dots, v_{n-1\bot}]}                  \\
                                                                                                 & \qquad \text{\cf \Cref{rule:semantics:grfrp:opsem:aux:intoVal:signal}}  \\
                  \CurrVal{ty_\sigma}{[v_{0\bot}, \dots, v_{n-1\bot}, v'_{n\bot}, v'_{n+1\bot}]} & = \CurrVal{ty_\sigma}{[v_{0\bot}, \dots, v_{n-1\bot}, \bot, \bot]}      \\
                                                                                                 & = \CurrVal{ty_\sigma}{[v_{0\bot}, \dots, v_{n-1\bot}]}                  \\
                                                                                                 & \qquad \text{\cf \Cref{rule:semantics:grfrp:opsem:aux:intoVal:signal}},
              \end{align*}
              which implies:
              \begin{align*}
                  s'_{n+2}                                                                       & = s_{n+1}                                            \\
                  \CurrVal{ty_\sigma}{[v_{0\bot}, \dots, v_{n-1\bot}, v'_{n\bot}, v'_{n+1\bot}]} & = \CurrVal{ty_\sigma}{[v_{0\bot}, \dots, v_{n\bot}]} \\
                  [(i'_{Tn+1}, t'_{Tn+1})^+]                                                     & = [(i_{Tn}, t_{Tn})^+].
              \end{align*}
    \end{itemize}
    We saw that in all cases we have:
    \begin{align*}
        s'_{n+2}                                                                       & = s_{n+1}                                            \\
        \CurrVal{ty_\sigma}{[v_{0\bot}, \dots, v_{n-1\bot}, v'_{n\bot}, v'_{n+1\bot}]} & = \CurrVal{ty_\sigma}{[v_{0\bot}, \dots, v_{n\bot}]} \\
        [(i'_{Tn+1}, t'_{Tn+1})^+]                                                     & = [(i_{Tn}, t_{Tn})^+],
    \end{align*}
    which concludes the proof for the \kwf{sample} operator.

    \paragraph{Component applications}
    Suppose $\sigma$ is $f(\sigma^+)$, the application of the component $f$ on the flow operations
    $\sigma^+$.

    The operational semantics rules \nameref{rule:semantics:grfrp:opsem:step:op:app:propag} and
    \nameref{rule:semantics:grfrp:opsem:step:op:app:value} ensure there exist states $s_{f0}$,
    \dots, $s_{fn+1}$, $s'_{fn+1}$, and $s'_{fn+2}$ corresponding to the component $f$, states
    $s_{i0}$, \dots, $s_{in+1}$, $s'_{in+1}$, and $s'_{in+2}$ for each input $\sigma$, bottom values
    $v_{i0\bot}$, \dots, $v_{in\bot}$, $v'_{in\bot}$, and $v'_{in+1\bot}$ and timer resets
    $[(i_{\sigma0}, t_{\sigma0})^+]$, \dots, $[(i_{\sigma n}, t_{\sigma n})^+]$,
    $[(i'_{\sigma n}, t'_{\sigma n})^+]$, and $[(i'_{\sigma n+1}, t'_{\sigma n+1})^+]$ for each
    input $\sigma$, and tuples of values $(v_{i0}^{\explast+})$, \dots, $(v_{in+1}^{\explast+})$,
    $(v_{in+1}^{\explast\prime+})$, $(v_{in+2}^{\explast\prime+})$, $(v_{o0}^{\explast+})$, \dots,
    $(v_{on+1}^{\explast+})$, $(v_{on+1}^{\explast\prime+})$, and $(v_{on+2}^{\explast\prime+})$
    such that:
    \begin{equation*}
        \mathfoot{\begin{array}{cc}
                s_0 = (s_{f0}, (v_{i0}^{\explast+}), (v_{o0}^{\explast+}), s_{i0}^+)                                   & \left(\sigmaOpSem{\gamma_0}{s_{i0}}{\sigma}{t_0}{s_{i1}}{v_{i0\bot}}{(i_{\sigma0}, t_{\sigma0})^+}\right)^+                             \\
                \dots                                                                                                  & \dots                                                                                                                                   \\
                s_{n+1} = (s_{fn+1}, (v_{in+1}^{\explast+}), (v_{on+1}^{\explast+}), s_{in+1}^+)                       & \left(\sigmaOpSem{\gamma_n}{s_{in}}{\sigma}{t_n}{s_{in+1}}{v_{in\bot}}{(i_{\sigma n}, t_{\sigma n})^+}\right)^+                         \\
                s'_{n+1} = (s'_{fn+1}, (v_{in+1}^{\explast\prime+}), (v_{on+1}^{\explast\prime+}), s_{in+1}^{\prime+}) & \left(\sigmaOpSem{\gamma'_n}{s_{in}}{\sigma}{t'_n}{s'_{in+1}}{v'_{in\bot}}{(i'_{\sigma n}, t'_{\sigma n})^+}\right)^+                   \\
                s'_{n+2} = (s'_{fn+2}, (v_{in+2}^{\explast\prime+}), (v_{on+2}^{\explast\prime+}), s_{in+2}^{\prime+}) & \left(\sigmaOpSem{\gamma'_{n+1}}{s'_{in+1}}{\sigma}{t'_{n+1}}{s'_{in+2}}{v'_{in+1\bot}}{(i'_{\sigma n+1}, t'_{\sigma n+1})^+}\right)^+.
            \end{array}}
    \end{equation*}

    The typing rule \nameref{rule:grust:typing:grfrp:flow_op:compapp} ensures there exists a tuple
    of types $(sty_i^+)$ such that:
    \begin{align*}
        \left(\flowTyped{\sigma}{sty_i}\right)^+.
    \end{align*}

    The reactive typing rule \nameref{rule:semantics:grfrp:react_typing:op:app} ensures that
    there exists a tuple of reactive types $(\rType_i^+)$ such that:
    \begin{equation*}
        \left(\opTypedReact{\Theta}{\sigma}{\rType_i}\right)^+ \qquad
        \compTypedReact{f}{\rType_i^+}{\rType^+}.
    \end{equation*}

    \Cref{eq:semantics:grfrp:theorems:determinism:deps:app} ensures that for all input $\sigma$ we
    have $\depSet{\sigma} \subseteq \depSet{f(\sigma^+)}$. It implies that for all input $\sigma$
    and $x$ in \depSet{\sigma} the current value of $x$ in the context without oversampling
    (encoded by \CurrVal{\tyCtx[x]}{[\gamma_0[x], \dots, \gamma_n[x]]}) equals its current value in the context with
    oversampling (encoded by \CurrVal{\tyCtx[x]}{[\gamma_0[x], \dots, \gamma_{n-1}[x], \gamma'_n[x], \gamma'_{n+1}[x]]}),
    and if $\Theta[x]$ equals $\dType$ then $\gamma'_n[x]$ is $\bot$.

    It is then possible to apply the structural induction hypothesis on each input $\sigma$, which
    implies that $s'_{in+2} = s_{in+2}$, $[(i'_{\sigma n}, t'_{\sigma n})^+]$ is empty, that the
    timer resets $[(i'_{\sigma n+1}, t'_{\sigma n+1})^+]$ equals $[(i_{\sigma n}, t_{\sigma n})^+]$,
    that the current value of the value of $\sigma$ without oversampling (encoded by the value \CurrVal{sty_i}{[v_{i0\bot}, \dots,
                v_{in\bot}]}) equals the one during the oversampling (encoded by the value \CurrVal{sty_i}{[v_{i0\bot}, \dots,
    v_{in-1\bot}, v'_{in\bot}, v'_{in+1\bot}]}), and that if $\rType_i = \dType$ then $v'_{in\bot}$
    is $\bot$.

    \begin{itemize}
        \item If $\left(v'_{in\bot} = \bot\right)^+$ (if the oversampling step has not produced any
              update), then the rule \nameref{rule:semantics:grfrp:opsem:step:op:app:propag} is
              applied. Therefore the state is unchanged ($s'_{n+1} = s_n$), all outputs $v'_{n\bot}$
              are $\bot$, and $[(i'_{Tn}, t'_{Tn})^+] = [((i'_{\sigma n},t'_{\sigma n})^+)^+]
                  \myeq{IH} \emptylist$.

              Furthermore, $\left(v'_{in\bot} = \bot\right)^+$ implies that the last value
              $\CurrVal{sty_i}{\mathfoot{[v_{i0\bot}, \dots, v_{in-1\bot}, v'_{in\bot}, v'_{in+1\bot}]}}$ equals
              $\CurrVal{sty_i}{\mathfoot{[v_{i0\bot}, \dots, v_{in-1\bot}, v'_{in+1\bot}]}}$, which equals
              $\CurrVal{sty_i}{\mathfoot{[v_{i0\bot}, \dots, v_{in-1\bot}, v_{in\bot}]}}$ by the induction hypothesis.
              \Cref{lem:semantics:grfrp:theorems:oversampling:CurrVal:last} ensures that $v'_{in+1\bot}$
              equals $v_{in\bot}$.

              In that case, because of the determinism of components
              (\cf \Cref{thm:semantics:grsync:theorems:determinism} for the determinism of the
              denotational semantics and \Cref{thm:semantics:grsync:theorems:equivalence} for the
              equivalence of operational and denotational semantics), both rules \nameref{rule:%
                  semantics:grfrp:opsem:step:op:app:value} and \nameref{rule:semantics:grfrp:opsem:%
                  step:op:app:propag} determine the output state $s'_{n+2}$, bottom values
              $(v_{n+1\bot}^{\prime+})$ and timer resets $[(i'_{Tn+1}, t'_{Tn+1})^+]$ of the
              component application purely based on the current state $s'_{n+1}$, input values
              $(v'_{in+1\bot})$, and input timer resets $[(i'_{\sigma n+1}, t'_{\sigma n+1})^+]$.
              Therefore, because $s'_{n+1} = s_n$, $v'_{in+1\bot} = v_{in\bot}$, and
              $[(i'_{\sigma n+1}, t'_{\sigma n+1})^+] \myeq{IH} [(i_{\sigma n}, t_{\sigma n})^+]$,
              we have:
              \begin{align*}
                  s'_{n+2}     & = s_{n+1}                                                                                                                                        \\
                  v'_{n+1\bot} & = v_{n\bot}                                                                                                                                      \\
                               & \Rightarrow \CurrVal{sty}{[v_{0\bot}, \dots, v_{n-1\bot}, v'_{n+1\bot}]} = \CurrVal{sty}{[v_{0\bot}, \dots, v_{n-1\bot}, v_{n\bot}]}             \\
                               & \qquad \text{by application of \Cref{lem:semantics:grfrp:theorems:oversampling:CurrVal:last}}                                                    \\
                               & \Rightarrow \CurrVal{sty}{[v_{0\bot}, \dots, v_{n-1\bot}, \bot, v'_{n+1\bot}]} = \CurrVal{sty}{[v_{0\bot}, \dots, v_{n-1\bot}, v_{n\bot}]}       \\
                               & \Rightarrow \CurrVal{sty}{[v_{0\bot}, \dots, v_{n-1\bot}, v'_{n\bot}, v'_{n+1\bot}]} = \CurrVal{sty}{[v_{0\bot}, \dots, v_{n-1\bot}, v_{n\bot}]} \\
                               & \qquad \text{\cf \Cref{rule:semantics:grfrp:opsem:aux:intoVal:signal,rule:semantics:grfrp:opsem:aux:intoVal:event}}                              \\
                               & [(i'_{Tn+1}, t'_{Tn+1})^+] = [((i'_{\sigma n+1}, t'_{\sigma n+1})^+)^+] \myeq{IH} [((i_{\sigma n}, t_{\sigma n})^+)^+] = [(i_{Tn}, t_{Tn})^+].
              \end{align*}
        \item If there exists $v'_{in\bot} \neq \bot$ (the oversampling has produce at least one
              input update), then the rule \nameref{rule:semantics:grfrp:opsem:step:op:app:value} is
              applied and we have:
              \begin{align*}
                   & (v'_{in} = \intoVal{sty_i}{v_{in}^{\explast}}{v'_{in\bot}})^+                                          \\
                   & \compOpsem{s_{fn}}{v_{in}^{\prime+}}{f}{s'_{fn+1}}{v_{on}^{\prime+}}                                   \\
                   & (v'_{n\bot} = \intoBot{sty}{v_{on}^{\explast}}{v'_{on}})^+                                             \\
                   & s'_{n+1} = (s'_{fn+1}, (v_{in+1}^{\explast\prime+}), (v_{on+1}^{\explast\prime+}), s_{in+1}^{\prime+})
                  = (s'_{fn+1}, v_{in}^{\prime+}, v_{on}^{\prime+}, s_{in+1}^{\prime+})                                     \\
                   & [(i'_{Tn}, t'_{Tn})^+] = [((i'_{\sigma n},t'_{\sigma n})^+)^+] \myeq{IH} \emptylist.
              \end{align*}

              For all input $\sigma$, if $\rType_i = \dType$ then $v'_{in\bot} \myeq{IH} \bot$.
              Then, if $\rType_i = \dType$, we have:
              \begin{equation*}
                  \begin{split}
                      v'_{in} & = \intoVal{sty_i}{v_{in}^{\explast}}{v'_{in\bot}}   \\
                              & \myeq{IH} \intoVal{sty_i}{v_{in}^{\explast}}{\bot}.
                  \end{split}
              \end{equation*}
              The application of the oversampling theorem of components
              (\Cref{thm:semantics:grsync:theorems:oversampling}) ensures that the state $s'_{fn+1}$
              remains equal to $s_{fn}$ and for all output value $v'_{on}$, if $\rType = \dType$
              then $v'_{on} = \intoVal{sty}{v_{on}^{\explast}}{\bot}$.

              Then, for all output value $v'_{on}$ such that $\rType = \dType$, we have:
              \begin{equation*}
                  \begin{split}
                      v'_{n\bot} & = \intoBot{sty}{v_{on}^{\explast}}{v'_{on}}                                                                                                                                                                     \\
                                 & = \intoBot{sty}{v_{on}^{\explast}}{\intoVal{sty}{v_{on}^{\explast}}{\bot}}                                                                                                                                      \\
                                 & = \bot                                                                                                                                                                                                          \\
                                 & \qquad \text{\cf \Cref{rule:semantics:grfrp:opsem:aux:intoVal:signal,rule:semantics:grfrp:opsem:aux:intoBot:signal,rule:semantics:grfrp:opsem:aux:intoVal:event,rule:semantics:grfrp:opsem:aux:intoBot:event}.}
                  \end{split}
              \end{equation*}

              In the case all $v_{in\bot}$ are equal to $\bot$ (and the rule
              \nameref{rule:semantics:grfrp:opsem:step:op:app:propag} is applied), we can
              still define $v_{in}$ to be $\intoVal{sty_i}{v_{in}^{\explast}}{v_{in\bot}}$.
              The oversampling theorem of components (\Cref{thm:semantics:grsync:theorems:oversampling})
              ensures that the resulting state of the step from $s_{fn}$ remains the same
              (\ie equal to $s_{fn+1}$ in the \nameref{rule:semantics:grfrp:opsem:step:op:app:propag}
              rule) and it defines output value $v_{on}$. In components, events are only
              used in \kwf{when} equations to emit other events or update the values of
              signals when they occur. Therefore, an absence of event does not produce any
              signal change or event occurrence, meaning that all output $v_{n\bot}$ are
              equal to \intoBot{sty}{v_{on}^{\explast}}{v_{on}} which is equal to $\bot$,
              \ie equal to the output value of the rule
              \nameref{rule:semantics:grfrp:opsem:step:op:app:propag}.
              This prove that the rule \nameref{rule:semantics:grfrp:opsem:step:op:app:value}
              and \nameref{rule:semantics:grfrp:opsem:step:op:app:propag} are equivalent,
              we can then suppose in the rest of this proof that the rule
              \nameref{rule:semantics:grfrp:opsem:step:op:app:value} is applied and that
              we have:
              \begin{align*}
                   & (v_{in} = \intoVal{sty_i}{v_{in}^{\explast}}{v_{in\bot}})^+                                                                   \\
                   & \compOpsem{s_{fn}}{v_{in}^+}{f}{s_{fn+1}}{v_{on}^+}                                                                           \\
                   & (v_{n\bot} = \intoBot{sty}{v_{on}^{\explast}}{v_{on}})^+                                                                      \\
                   & s_{n+1} = (s_{fn+1}, (v_{in+1}^{\explast+}), (v_{on+1}^{\explast+}), s_{in+1}^+) = (s_{fn+1}, v_{in}^+, v_{on}^+, s_{in+1}^+) \\
                   & [(i_{Tn}, t_{Tn})^+] = [((i_{\sigma n}, t_{\sigma n})^+)^+] = [(i'_{Tn+1}, t'_{Tn+1})^+].
              \end{align*}

              Similarly, we suppose in the rest of this proof that the rule
              \nameref{rule:semantics:grfrp:opsem:step:op:app:value} is applied for the oversampled
              simulation and that we have:
              \begin{align*}
                   & (v'_{in+1} = \intoVal{sty_i}{v_{in+1}^{\explast\prime}}{v'_{in+1\bot}} = \intoVal{sty_i}{v'_{in}}{v'_{in+1\bot}})^+                                                              \\
                   & \compOpsem{s_{fn}}{v_{in+1}^{\prime+}}{f}{s'_{fn+2}}{v_{on+1}^{\prime+}}                                                                                                         \\
                   & (v'_{n+1\bot} = \intoBot{sty}{v_{on+1}^{\explast\prime}}{v'_{on+1}} = \intoBot{sty}{v'_{on}}{v'_{on+1}})^+                                                                       \\
                   & s'_{n+2} = (s'_{fn+2}, (v_{in+2}^{\explast\prime+}), (v_{on+2}^{\explast\prime+}), s_{in+2}^{\prime+}) = (s'_{fn+2}, v_{in+1}^{\prime+}, v_{on+1}^{\prime+}, s_{in+2}^{\prime+}) \\
                   & [(i'_{Tn+1}, t'_{Tn+1})^+] = [((i'_{\sigma n+1}, t'_{\sigma n+1})^+)^+] \myeq{IH} [((i_{\sigma n}, t_{\sigma n})^+)^+].
              \end{align*}

              For all input $v'_{in+1}$, we have three cases:
              \begin{enumerate}
                  \item Suppose $\rType_i = \dType$. It ensures that $v'_{in\bot} \myeq{IH} \bot$,
                        which in turn implies:
                        \begin{align*}
                            \CurrVal{sty_i}{[v_{i0\bot}, \dots, v_{in\bot}]}               & \myeq{IH} \CurrVal{sty_i}{[v_{i0\bot}, \dots, v_{in-1\bot}, v'_{in\bot}, v'_{in+1\bot}]}                            \\
                            \CurrVal{sty_i}{[v_{i0\bot}, \dots, v_{in-1\bot}, v_{in\bot}]} & = \CurrVal{sty_i}{[v_{i0\bot}, \dots, v_{in-1\bot}, \bot, v'_{in+1\bot}]}                                           \\
                                                                                           & = \CurrVal{sty_i}{[v_{i0\bot}, \dots, v_{in-1\bot}, v'_{in+1\bot}]}                                                 \\
                                                                                           & \qquad \text{\cf \Cref{rule:semantics:grfrp:opsem:aux:intoVal:signal,rule:semantics:grfrp:opsem:aux:intoVal:event}} \\
                                                                                           & \Rightarrow \quad v_{in\bot} = v'_{in+1\bot}                                                                        \\
                                                                                           & \qquad \text{\cf \Cref{lem:semantics:grfrp:theorems:oversampling:CurrVal:last}},
                        \end{align*}
                        then, if there exist $ty_i$ such that $sty_i = ty_i$ (\ie is a signal type):
                        \begin{align*}
                            v'_{in+1} & = \intoVal{ty_i}{v'_{in}}{v'_{in+1\bot}}                                     \\
                                      & = \intoVal{ty_i}{\intoVal{ty_i}{v_{in}^{\explast}}{v'_{in\bot}}}{v_{in\bot}} \\
                                      & = \intoVal{ty_i}{\intoVal{ty_i}{v_{in}^{\explast}}{\bot}}{v_{in\bot}}        \\
                                      & = \intoVal{ty_i}{v_{in}^{\explast}}{v_{in\bot}}                              \\
                                      & \qquad \text{\cf \Cref{rule:semantics:grfrp:opsem:aux:intoVal:signal}}       \\
                                      & = v_{in},
                        \end{align*}
                        and, if there exist $ty_i$ such that $sty_i = \mayB{ty_i}$ (\ie is an event type):
                        \begin{align*}
                            v'_{in+1} & = \intoVal{\mayB{ty_i}}{v'_{in}}{v'_{in+1\bot}}                       \\
                                      & = \intoVal{\mayB{ty_i}}{\_}{v_{in\bot}}                               \\
                                      & = \intoVal{\mayB{ty_i}}{v_{in}^{\explast}}{v_{in\bot}}                \\
                                      & \qquad \text{\cf \Cref{rule:semantics:grfrp:opsem:aux:intoVal:event}} \\
                                      & = v_{in}.
                        \end{align*}
                  \item Suppose $\rType_i \neq \dType$ and $v'_{in+1\bot} = \bot$. Then, because
                        only signals can be continuous in the reactive type system, there exist
                        $ty_i$ such that $sty_i = ty_i$. It implies:
                        \begin{align*}
                            \CurrVal{ty_i}{[v_{i0\bot}, \dots, v_{in-1\bot}, v_{in\bot}]} & \myeq{IH} \CurrVal{ty_i}{[v_{i0\bot}, \dots, v_{in-1\bot}, v'_{in\bot}, v'_{in+1\bot}]} \\
                                                                                          & = \CurrVal{ty_i}{[v_{i0\bot}, \dots, v_{in-1\bot}, v'_{in\bot}, \bot]}                  \\
                                                                                          & = \CurrVal{ty_i}{[v_{i0\bot}, \dots, v_{in-1\bot}, v'_{in\bot}]}                        \\
                                                                                          & \qquad \text{\cf \Cref{rule:semantics:grfrp:opsem:aux:intoVal:signal}}                  \\
                                                                                          & \Rightarrow \quad v_{in\bot} = v'_{in\bot}                                              \\
                                                                                          & \qquad \text{\cf \Cref{lem:semantics:grfrp:theorems:oversampling:CurrVal:last}},
                        \end{align*}
                        then:
                        \begin{align*}
                            v'_{in+1} & = \intoVal{ty_i}{v'_{in}}{v'_{in+1\bot}}                               \\
                                      & = \intoVal{ty_i}{v'_{in}}{\bot}                                        \\
                                      & = \intoVal{ty_i}{\intoVal{ty_i}{v_{in}^{\explast}}{v'_{in\bot}}}{\bot} \\
                                      & = \intoVal{ty_i}{\intoVal{ty_i}{v_{in}^{\explast}}{v_{in\bot}}}{\bot}  \\
                                      & = \intoVal{ty_i}{v_{in}^{\explast}}{v_{in\bot}}                        \\
                                      & \qquad \text{\cf \Cref{rule:semantics:grfrp:opsem:aux:intoVal:signal}} \\
                                      & = v_{in}.
                        \end{align*}
                  \item Suppose $\rType_i \neq \dType$ and $v'_{in+1\bot} \neq \bot$. Similarly,
                        because only signals can be continuous in the reactive type system, there
                        exist $ty_i$ such that $sty_i = ty_i$. It implies:
                        \begin{align*}
                            \CurrVal{sty_i}{[v_{i0\bot}, \dots, v_{in-1\bot}, v_{in\bot}]} & \myeq{IH} \CurrVal{sty_i}{[v_{i0\bot}, \dots, v_{in-1\bot}, v'_{in\bot}, v'_{in+1\bot}]} \\
                                                                                           & = v'_{in+1\bot},
                        \end{align*}
                        \begin{align*}
                            v_{in} & = \intoVal{sty_i}{v_{in}^{\explast}}{v_{in\bot}}                                                     \\
                                   & = \intoVal{sty_i}{\intoVal{sty_i}{v_{in-1}^{\explast}}{v_{in-1\bot}}}{v_{in\bot}}                    \\
                                   & \qquad \text{\cf semantics rules}                                                                    \\
                                   & = \intoVal{sty_i}{v_{in-1}^{\explast}}{\CurrVal{sty_i}{[v_{in-1\bot}, v_{in\bot}]}}                  \\
                                   & \qquad \text{\cf \Cref{rule:semantics:grfrp:opsem:aux:intoVal:signal}}                               \\
                                   & \dots                                                                                                \\
                                   & = \intoVal{sty_i}{v_{i0}^{\explast}}{\CurrVal{sty_i}{[v_{i0\bot}, \dots, v_{in-1\bot}, v_{in\bot}]}} \\
                                   & = \intoVal{sty_i}{v_{i0}^{\explast}}{v'_{in+1\bot}}                                                  \\
                                   & = v'_{in+1\bot}                                                                                      \\
                                   & = \intoVal{sty_i}{v'_{in}}{v'_{in+1\bot}}                                                            \\
                                   & = v'_{in+1}.
                        \end{align*}
              \end{enumerate}
              In all cases, we have $v'_{in+1} = v_{in}$.

              Because components are deterministic (\cf \Cref{thm:semantics:grsync:theorems:determinism}
              for the determinism of the denotational semantics and \Cref{thm:semantics:grsync:theorems:equivalence}
              for the equivalence of operational and denotational semantics), we have:
              \begin{align*}
                  s'_{fn+2} & = s_{fn+1} & (v'_{on+1} & = v_{on})^+,
              \end{align*}
              which ensures that for all output $v'_{n+1\bot}$ we have:
              \begin{align*}
                  v'_{n+1\bot} & = \intoBot{sty}{v'_{on}}{v'_{on+1}}                                                                        \\
                  v'_{n+1\bot} & = \intoBot{sty}{v'_{on}}{v_{on}}                                                                           \\
                  v'_{n+1\bot} & = \intoBot{sty}{\intoVal{sty}{v'_{on-1}}{v'_{n\bot}}}{v_{on}}                                              \\
                               & \dots                                                                                                      \\
                  v'_{n+1\bot} & = \intoBot{sty}{\intoVal{sty}{v_{o0}}{\CurrVal{sty}{[v_{0\bot}, \dots, v_{n-1\bot}, v'_{n\bot}]}}}{v_{on}} \\
                  v'_{n+1\bot} & = \left\{
                  \begin{array}{ll}
                      v_{on} & \text{if} \, sty = ty \wedge v_{on} \neq \intoVal{sty}{v_{o0}}{\CurrVal{sty}{[v_{0\bot}, \dots, v_{n-1\bot}, v'_{n\bot}]}} \\
                      \bot   & \text{if} \, sty = ty \wedge v_{on} =    \intoVal{sty}{v_{o0}}{\CurrVal{sty}{[v_{0\bot}, \dots, v_{n-1\bot}, v'_{n\bot}]}} \\
                      v      & \text{if} \, sty = \mayB{ty} \wedge v_{on} = \some{v}                                                                      \\
                      \bot   & \text{if} \, sty = \mayB{ty} \wedge v_{on} = \none,
                  \end{array}
                  \right.                                                                                                                   \\
                  v_{n\bot}    & = \left\{
                  \begin{array}{ll}
                      v_{on} & \text{if} \, sty = ty \wedge v_{on} \neq \intoVal{sty}{v_{o0}}{\CurrVal{sty}{[v_{0\bot}, \dots, v_{n-1\bot}]}} \\
                      \bot   & \text{if} \, sty = ty \wedge v_{on} =    \intoVal{sty}{v_{o0}}{\CurrVal{sty}{[v_{0\bot}, \dots, v_{n-1\bot}]}} \\
                      v      & \text{if} \, sty = \mayB{ty} \wedge v_{on} = \some{v}                                                          \\
                      \bot   & \text{if} \, sty = \mayB{ty} \wedge v_{on} = \none,
                  \end{array}
                  \right.
              \end{align*}
              then, if $sty = \mayB{ty}$ we have $v'_{n+1\bot} = v_{n\bot}$ and thus (knowing that
              $\rType = \dType$ so that $v'_{n\bot} = \bot$)
              $\CurrVal{sty}{[v_{0\bot}, \dots, v_{n-1\bot}, v'_{n\bot}, v'_{n+1\bot}]} =
                  \CurrVal{sty}{[v_{0\bot}, \dots, v_{n-1\bot}, v_{n\bot}]}$, and if $sty = ty$ we have:
              \begin{multline*}
                  \CurrVal{sty}{[v_{0\bot}, \dots, v_{n-1\bot}, v'_{n\bot}, v'_{n+1\bot}]} \\
                  = \left\{
                  \begin{array}{l}
                      v_{on}                                                                                                            \\
                      \qquad  \text{if} \, v_{on} \neq \intoVal{ty}{v_{o0}}{\CurrVal{sty}{[v_{0\bot}, \dots, v_{n-1\bot}, v'_{n\bot}]}} \\
                      \CurrVal{sty}{[v_{0\bot}, \dots, v_{n-1\bot}, v'_{n\bot}]}                                                        \\
                      \qquad  \text{if} \, v_{on} = \intoVal{ty}{v_{o0}}{\CurrVal{sty}{[v_{0\bot}, \dots, v_{n-1\bot}, v'_{n\bot}]}}
                  \end{array}
                  \right.                                                                        \\
                  = \left\{
                  \begin{array}{ll}
                      v_{on} & \text{if} \, \CurrVal{sty}{[v_{0\bot}, \dots, v_{n-1\bot}, v'_{n\bot}]} = \bot \wedge v_{on} \neq v_{o0} \\
                      v_{on} & \text{if} \, \CurrVal{sty}{[v_{0\bot}, \dots, v_{n-1\bot}, v'_{n\bot}]} = v \wedge v_{on} \neq v         \\
                      \bot   & \text{if} \, \CurrVal{sty}{[v_{0\bot}, \dots, v_{n-1\bot}, v'_{n\bot}]} = \bot \wedge v_{on} = v_{o0}    \\
                      v      & \text{if} \, \CurrVal{sty}{[v_{0\bot}, \dots, v_{n-1\bot}, v'_{n\bot}]} = v \wedge v_{on} = v
                  \end{array}
                  \right.                                                                        \\
                  = \left\{
                  \begin{array}{ll}
                      v_{on} & \text{if} \, \CurrVal{sty}{[v_{0\bot}, \dots, v_{n-1\bot}, v'_{n\bot}]} = \bot \wedge v_{on} \neq v_{o0} \\
                      \bot   & \text{if} \, \CurrVal{sty}{[v_{0\bot}, \dots, v_{n-1\bot}, v'_{n\bot}]} = \bot \wedge v_{on} = v_{o0}    \\
                      v_{on} & \text{if} \, \CurrVal{sty}{[v_{0\bot}, \dots, v_{n-1\bot}, v'_{n\bot}]} \neq \bot,
                  \end{array}
                  \right.
              \end{multline*}
              similarly we have:
              \begin{multline*}
                  \CurrVal{sty}{[v_{0\bot}, \dots, v_{n-1\bot}, v_{n\bot}]} \\
                  = \left\{
                  \begin{array}{l}
                      v_{on}                                                                                                \\
                      \qquad  \text{if} \, v_{on} \neq \intoVal{ty}{v_{o0}}{\CurrVal{sty}{[v_{0\bot}, \dots, v_{n-1\bot}]}} \\
                      \CurrVal{sty}{[v_{0\bot}, \dots, v_{n-1\bot}]}                                                        \\
                      \qquad  \text{if} \, v_{on} = \intoVal{ty}{v_{o0}}{\CurrVal{sty}{[v_{0\bot}, \dots, v_{n-1\bot}]}}
                  \end{array}
                  \right.                                                                        \\
                  = \left\{
                  \begin{array}{ll}
                      v_{on} & \text{if} \, \CurrVal{sty}{[v_{0\bot}, \dots, v_{n-1\bot}]} = \bot \wedge v_{on} \neq v_{o0} \\
                      v_{on} & \text{if} \, \CurrVal{sty}{[v_{0\bot}, \dots, v_{n-1\bot}]} = v \wedge v_{on} \neq v         \\
                      \bot   & \text{if} \, \CurrVal{sty}{[v_{0\bot}, \dots, v_{n-1\bot}]} = \bot \wedge v_{on} = v_{o0}    \\
                      v      & \text{if} \, \CurrVal{sty}{[v_{0\bot}, \dots, v_{n-1\bot}]} = v \wedge v_{on} = v
                  \end{array}
                  \right.                                                                        \\
                  = \left\{
                  \begin{array}{ll}
                      v_{on} & \text{if} \, \CurrVal{sty}{[v_{0\bot}, \dots, v_{n-1\bot}]} = \bot \wedge v_{on} \neq v_{o0} \\
                      \bot   & \text{if} \, \CurrVal{sty}{[v_{0\bot}, \dots, v_{n-1\bot}]} = \bot \wedge v_{on} = v_{o0}    \\
                      v_{on} & \text{if} \, \CurrVal{sty}{[v_{0\bot}, \dots, v_{n-1\bot}]} \neq \bot,
                  \end{array}
                  \right.
              \end{multline*}
              so that we have $\CurrVal{sty}{[v_{0\bot}, \dots, v_{n-1\bot}, v'_{n\bot}, v'_{n+1\bot}]} =
                  \CurrVal{sty}{[v_{0\bot}, \dots, v_{n-1\bot}, v_{n\bot}]}$.

              For the new state $s'_{n+2}$, we have:
              \begin{equation*}
                  \begin{split}
                      s'_{n+2} & = (s'_{fn+2}, v_{in+1}^{\prime+}, v_{on+1}^{\prime+}, s_{in+2}^{\prime+}) \\
                               & = (s_{fn+1}, v_{in}^+, v_{on}^+, s_{in+1}^+)                              \\
                               & = s_{n+1}.
                  \end{split}
              \end{equation*}
    \end{itemize}
    In all cases, the state following the oversampling is equal to the not oversampled simulation
    ($s'_{n+2} = s_{n+1}$), the current values of all output flows with and without oversampling are
    the same ($\CurrVal{sty}{[v_{0\bot}, \dots, v_{n-1\bot}, v'_{n\bot}, v'_{n+1\bot}]} =
        \CurrVal{sty}{[v_{0\bot}, \dots, v_{n-1\bot}, v_{n\bot}]}$), we have no resets when oversampling
    ($[(i'_{Tn}, t'_{Tn})^+]$ is empty), the timer resets following the oversampling are equal to
    the ones without oversampling ($[(i'_{Tn+1}, t'_{Tn+1})^+] = [(i_{Tn}, t_{Tn})^+]$), and for all
    discrete output flow (\ie such that $\rType = \dType$) the oversampling induce no changes
    ($v'_{n\bot} = \bot$).

    \paragraph{Other flow operations}

    The proof of other flow operators follows the same principle as the proof of the
    \sample{\sigma}{T} operation:

    \begin{enumerate}
        \item The different operational semantics rules for the flow operation ensure the existence
              of sequences of states corresponding to the inner flow operations ($s_{\sigma 0}$,
              \dots, $s_{\sigma n+1}$, $s'_{\sigma n+1}$ and $s'_{\sigma n+2}$ corresponding to
              $\sigma$ for the \sample{\sigma}{T} operation), sequence of bottom values
              corresponding to the updates of those inner flow operation ($v_{\sigma 0\bot}$, \dots,
              $v_{\sigma n\bot}$, $v'_{\sigma n\bot}$ and $v'_{\sigma n+1\bot}$ for the
              \sample{\sigma}{T} operation), along with their sequence of lists of timer resets
              ($[(i_{\sigma0},t_{\sigma0})^+]$, \dots, $[(i_{\sigma n}, t_{\sigma n})^+]$,
              $[(i'_{\sigma n}, t'_{\sigma n})^+]$ and $[(i'_{\sigma n+1}, t'_{\sigma n+1})^+]$ for
              the \sample{\sigma}{T} operation), and sequence of state values ($(v_{m0\bot})$,
              \dots, $(v_{mn+1\bot})$, $(v'_{mn+1\bot})$, $(v'_{mn+2\bot})$ for the
              \sample{\sigma}{T} operation).
        \item The typing rules, reactive typing rules, and dependency sets gives makes it possible
              to use the induction hypothesis on the inner flow operations. This adds constraints on
              the update values, timer resets, and states of inner flow operations.
        \item We process the oversampling step first (at $t'_n$) to prove that there is no timer
              resets produced at that time and no updates for discrete flows. For flow operations
              depending on timers, we know that the presence of a timer identifier ($i_T$) in an
              update context ($\gamma$) can only come from the propagation of a timer message
              (\timerMess{i_T}), as it is only shown by the rule
              \nameref{rule:semantics:grfrp:opsem:step:service:timer}. Therefore, timer identifiers
              are not present in update contexts during the oversampling. It restrains the possible
              step transitions for the oversampling (only the rules
              \nameref{rule:semantics:grfrp:opsem:step:op:sample:propag} and
              \nameref{rule:semantics:grfrp:opsem:step:op:sample:value} are suitable for the
              \sample{\sigma}{T}, but as the induction hypothesis and the reactive type system
              ensures $v'_{\sigma n\bot} = \bot$, the rule
              \nameref{rule:semantics:grfrp:opsem:step:op:sample:value} is disabled).
        \item We then prove that, in any transition case from the operational semantics, the step
              after oversampling is equivalent to the step without oversampling.
              For example, we proved for the \sample{\sigma}{T} flow operation that in all cases we
              have:
              \begin{align*}
                  s'_{n+2}                                                                 & = s_{n+1}                                      \\
                  \CurrVal{sty}{[v_{0\bot}, \dots, v_{n-1\bot}, v'_{n\bot}, v'_{n+1\bot}]} & = \CurrVal{sty}{[v_{0\bot}, \dots, v_{n\bot}]} \\
                  [(i'_{Tn+1}, t'_{Tn+1})^+]                                               & = [(i_{Tn}, t_{Tn})^+].
              \end{align*}
    \end{enumerate}

    \setcounter{paragraph}{0}
\end{skippableproof}

\begin{lemma}[Oversampling of \GRFRP reactive statements]
    \label{lem:semantics:grfrp:theorems:oversampling:stmts}
    Let $stmt^+$ be a sequence of statements in a program \G with a typing context \tyCtx, and
    let $\Theta$ be a reactive typing context such that:
    \begin{equation*}
        \left(\stmtTyped{stmt}\right)^* \qquad \stmtTypedReact{\Theta}{stmt^+}.
    \end{equation*}
    For
    all $n \ge 1$;
    all states $s_0$, \dots, $s_{n+1}$, $s'_{n+1}$, and $s'_{n+2}$ corresponding to $stmt^+$;
    all input update contexts $\gamma_{i0}$, \dots, $\gamma_{in}$, $\gamma'_{in}$, and $\gamma'_{in+1}$;
    all output update contexts $\gamma_{o0}$, \dots, $\gamma_{on}$, $\gamma'_{on}$, and $\gamma'_{on+1}$;
    all triggering inputs $(i_0, t_0, v_0)$, \dots, $(i_n, t_n, v_n)$, $(i'_n, t'_n, v'_n)$, and $(i'_{n+1}, t'_{n+1}, v'_{n+1})$;
    all lists of outputs $[(x_0, v_0)^+]$, \dots, $[(x_n, v_n)^+]$, $[(x'_n, v'_n)^+]$, and $[(x'_{n+1}, v'_{n+1})^+]$; and
    all lists of timer resets $[(i_{T0}, t_{T0})^+]$, \dots, $[(i_{Tn}, t_{Tn})^+]$, $[(i'_{Tn}, t'_{Tn})^+]$, and $[(i'_{Tn+1}, t'_{Tn+1})^+]$
    such that:
    \begin{equation*}
        \begin{array}{c}
            t'_n \le t'_{n+1} \qquad (i'_{n+1}, t'_{n+1}, v'_{n+1}) = (i_n, t_n, v_n) \qquad  s_0 = \initstate{stmt^+}                                       \\
            \forall k \in \{0, \dots, n\}, \quad \stmtOpSem{\gamma_{ik}}{s_k}{stmt^+}{i_k, t_k, v_k}{\gamma_{ok}}{s_{k+1}}{(x_k, v_k)^+}{(i_{Tk}, t_{Tk})^+} \\
            \stmtOpSem{\gamma'_{in}}{s_n}{stmt^+}{i'_n, t'_n, v'_n}{\gamma'_{on}}{s'_{n+1}}{(x'_n, v'_n)^+}{(i'_{Tn}, t'_{Tn})^+}                            \\
            \stmtOpSem{\gamma'_{in+1}}{s'_{n+1}}{stmt^+}{i'_{n+1}, t'_{n+1}, v'_{n+1}}{\gamma'_{on+1}}{s'_{n+2}}{(x'_{n+1}, v'_{n+1})^+}{(i'_{Tn+1}, t'_{Tn+1})^+},
        \end{array}
    \end{equation*}
    if $i'_n$ is not a timer identifier and not in \importSet{stmt^+}, then we have: % \todo{timer identifier set}, best would be not a timer nor a dependency of the statements
    \begin{multline*}
        \left[\begin{array}{l}
                \forall x \in \depSet{stmt^+}, \; \left(\Theta[x] = \dType \Rightarrow \gamma'_{in}[x] = \bot\right) \\
                \quad \wedge \left(\CurrVal{\tyCtx[x]}{[\gamma_{i0}[x], \dots, \gamma_{in}[x]]} = \CurrVal{\tyCtx[x]}{[\gamma_{i0}[x], \dots, \gamma_{in-1}[x], \gamma'_{in}[x], \gamma'_{in+1}[x]]}\right)
            \end{array}\right]
        \\ \Rightarrow \\
        \begin{array}{l}
            \left[\begin{array}{l}
                          \forall y \in \declSet{stmt^+}, \; \left(\Theta[y] = \dType \Rightarrow \gamma'_{on}[y] = \bot\right) \\
                          \quad \wedge \left(\CurrVal{\tyCtx[y]}{[\gamma_{o0}[y], \dots, \gamma_{on}[y]]} = \CurrVal{\tyCtx[y]}{[\gamma_{o0}[y], \dots, \gamma_{on-1}[y], \gamma'_{on}[y], \gamma'_{on+1}[y]]}\right)
                      \end{array}\right] \\
            \wedge \quad s'_{n+2} = s_{n+1} \wedge [(x'_{n+1}, v'_{n+1})^+] = [(x_n, v_n)^+] \wedge [(i'_{Tn+1}, t'_{Tn+1})^+] = [(i_{Tn}, t_{Tn})^+]                                                                                                                                                                                                                                        \\
            \wedge \quad [(x'_n, v'_n)^+] = \emptylist \wedge [(i'_{Tn}, t'_{Tn})^+] = \emptylist.
        \end{array}
    \end{multline*}

    This lemma formalizes the behavior of flow statements, $stmt^+$, during an oversampling step.
    It states that if, after $n$ simulation steps, all discrete dependencies of $stmt^+$ (of type
    \dType) are oversampled at step $n+1$, then this intermediate step has no lasting effect.
    Specifically, a computation performed from this new state yields the same result as one
    performed from the state at step $n$. This implies that the internal state does not change
    ($s'_{n+2} = s_{n+1}$), timer resets are preserved ($[(i'_{Tn}, t'_{Tn})^+] = \emptylist$ and
    $[(i'_{Tn+1}, t'_{Tn+1})^+] = [(i_{Tn}, t_{Tn})^+]$), and discrete flows are oversampled
    ($\gamma'_{on}[y] = \bot$). Furthermore, exported values are unchanged
    ($[(x'_n, v'_n)^+] = \emptylist$ and $[(x'_{n+1}, v'_{n+1})^+] = [(x_n, v_n)^+]$), and flows
    have identical current values ($\CurrVal{\tyCtx[y]}{[\dots, \gamma_{on}[y]]} =
        \CurrVal{\tyCtx[y]}{[\dots, \gamma_{on-1}[y], \gamma'_{on}[y], \gamma'_{on+1}[y]]}$).
\end{lemma}
\begin{skippableproof}
    Let $stmt^+$ be a sequence of statements in a program \G with a typing context \tyCtx, and
    let $\Theta$ be a reactive typing context such that:
    \begin{equation*}
        \left(\stmtTyped{stmt}\right)^* \qquad \stmtTypedReact{\Theta}{stmt^+}.
    \end{equation*}
    Let $n \ge 1$;
    let $s_0$, \dots, $s_{n+1}$, $s'_{n+1}$, and $s'_{n+2}$ be states corresponding to $stmt^+$;
    let $\gamma_0$, \dots, $\gamma_n$, $\gamma'_n$, and $\gamma'_{n+1}$ be input update contexts;
    let $\gamma_{o0}$, \dots, $\gamma_{on}$, $\gamma'_{on}$, and $\gamma'_{on+1}$ be output update contexts;
    let $(i_0, t_0, v_0)$, \dots, $(i_n, t_n, v_n)$, $(i'_n, t'_n, v'_n)$, and $(i'_{n+1}, t'_{n+1}, v'_{n+1})$ be triggering inputs;
    let $[(x_0, v_0)^+]$, \dots, $[(x_n, v_n)^+]$, $[(x'_n, v'_n)^+]$, and $[(x'_{n+1}, v'_{n+1})^+]$ be lists of outputs; and
    let $[(i_{T0}, t_{T0})^+]$, \dots, $[(i_{Tn}, t_{Tn})^+]$, $[(i'_{Tn}, t'_{Tn})^+]$, and $[(i'_{Tn+1}, t'_{Tn+1})^+]$ be lists of timer resets
    such that:
    \begin{equation*}
        \begin{array}{c}
            t'_n \le t'_{n+1} \qquad (i'_{n+1}, t'_{n+1}, v'_{n+1}) = (i_n, t_n, v_n) \qquad  s_0 = \initstate{stmt^+}                                       \\
            \forall k \in \{0, \dots, n\}, \quad \stmtOpSem{\gamma_{ik}}{s_k}{stmt^+}{i_k, t_k, v_k}{\gamma_{ok}}{s_{k+1}}{(x_k, v_k)^+}{(i_{Tk}, t_{Tk})^+} \\
            \stmtOpSem{\gamma'_n}{s_n}{stmt^+}{i'_n, t'_n, v'_n}{\gamma'_{on}}{s'_{n+1}}{(x'_n, v'_n)^+}{(i'_{Tn}, t'_{Tn})^+}                               \\
            \stmtOpSem{\gamma'_{n+1}}{s'_{n+1}}{stmt^+}{i'_{n+1}, t'_{n+1}, v'_{n+1}}{\gamma'_{on+1}}{s'_{n+2}}{(x'_{n+1}, v'_{n+1})^+}{(i'_{Tn+1}, t'_{Tn+1})^+}.
        \end{array}
    \end{equation*}
    Suppose $i'_n$ is not a timer identifier and not in \importSet{stmt^+}, and that
    we have:
    \begin{multline*}
        \forall x \in \depSet{stmt^+}, \; \left(\Theta[x] = \dType \Rightarrow \gamma'_n[x] = \bot\right) \\
        \wedge \left(\CurrVal{\tyCtx[x]}{[\gamma_0[x], \dots, \gamma_n[x]]} = \CurrVal{\tyCtx[x]}{[\gamma_0[x], \dots, \gamma_{n-1}[x], \gamma'_n[x], \gamma'_{n+1}[x]]}\right).
    \end{multline*}

    We proceed by structural induction on $stmt^+$:.
    \begin{itemize}
        \item Suppose $stmt^+$ is the empty block of statements $\{\skp\}$.

              The operational rule \nameref{rule:semantics:grfrp:opsem:step:stmt:skip} ensures that:
              \begin{align*}
                  s_0      & = () & \gamma_{o0}    & = \gamma_0      & s_1      & = () & (x_0, v_0)^+           & = \emptylist & (i_{T0}, t_{T0})^+       & = \emptylist  \\
                  s_n      & = () & \gamma_{on}    & = \gamma_n      & s_{n+1}  & = () & (x_n, v_n)^+           & = \emptylist & (i_{Tn}, t_{Tn})^+       & = \emptylist  \\
                  s_n      & = () & \gamma'_{on}   & = \gamma'_n     & s'_{n+1} & = () & (x'_n, v'_n)^+         & = \emptylist & (i'_{Tn}, t'_{Tn})^+     & = \emptylist  \\
                  s'_{n+1} & = () & \gamma'_{on+1} & = \gamma'_{n+1} & s'_{n+2} & = () & (x'_{n+1}, v'_{n+1})^+ & = \emptylist & (i'_{Tn+1}, t'_{Tn+1})^+ & = \emptylist.
              \end{align*}

              Therefore we have:
              \begin{align*}
                  [(i'_{Tn}, t'_{Tn})^+]     & = \emptylist           &
                                             &                        &
                  [(x'_n, v'_n)^+]           & = \emptylist             \\
                  [(i'_{Tn+1}, t'_{Tn+1})^+] & = [(i_{Tn}, t_{Tn})^+] &
                  s'_{n+2}                   & = s_{n+1}              &
                  [(x'_{n+1}, v'_{n+1})^+]   & = [(x_n, v_n)^+].
              \end{align*}

              \Cref{eq:semantics:grfrp:theorems:determinism:sets:decl:skip} states that
              $\declSet{\skp} = \emptyset$, which implies that:
              \begin{equation*}
                  \begin{array}{l}
                      \forall y \in \declSet{\skp}, \; \left(\Theta[y] = \dType \Rightarrow \gamma'_{on}[y] = \bot\right) \\
                      \quad \wedge \left(\CurrVal{\tyCtx[y]}{[\gamma_{o0}[y], \dots, \gamma_{on}[y]]} = \CurrVal{\tyCtx[y]}{[\gamma_{o0}[y], \dots, \gamma_{on-1}[y], \gamma'_{on}[y], \gamma'_{on+1}[y]]}\right).
                  \end{array}
              \end{equation*}

        \item Suppose $stmt^+$ can be decomposed into $\import{x}, stmt^+$.

              The operational rules \nameref{rule:semantics:grfrp:opsem:step:stmt:import:present}
              and \nameref{rule:semantics:grfrp:opsem:step:stmt:import:absent} ensure that there
              exist update environments $\gamma_{i0}$, \dots, $\gamma_{in}$, $\gamma'_{in}$, and
              $\gamma'_{in+1}$ such that for all $k \in \{0, \dots, n\}$:
              \begin{equation*}
                  \begin{array}{c}
                      \gamma_{ik} =
                      \left\{
                      \begin{array}{ll}
                          \gamma_k \ctxconcat [v_k/x] & \text{if} \; x = i_k    \\
                          \gamma_k                    & \text{if} \; x \neq i_k
                      \end{array}
                      \right. \\
                      \stmtOpSem{\gamma_{ik}}{s_k}{stmt^+}{i_k, t_k, v_k}{\gamma_{ok}}{s_{k+1}}{(x_k, v_k)^+}{(i_{Tk}, t_{Tk})^+}
                  \end{array}
              \end{equation*}
              and
              \begin{equation*}
                  \begin{array}{c}
                      \gamma'_{in} =
                      \left\{
                      \begin{array}{ll}
                          \gamma'_n \ctxconcat [v'_n/x] & \text{if} \; x = i'_n    \\
                          \gamma'_n                     & \text{if} \; x \neq i'_n
                      \end{array}
                      \right.                                                                                                               \\
                      \stmtOpSem{\gamma'_{in}}{s_n}{stmt^+}{i'_n, t'_n, v'_n}{\gamma'_{on}}{s'_{n+1}}{(x'_n, v'_n)^+}{(i'_{Tn}, t'_{Tn})^+} \\
                      \gamma'_{in+1} =
                      \left\{
                      \begin{array}{ll}
                          \gamma'_{n+1} \ctxconcat [v'_{n+1}/x] & \text{if} \; x = i'_{n+1}    \\
                          \gamma'_{n+1}                         & \text{if} \; x \neq i'_{n+1}
                      \end{array}
                      \right.                                                                                                               \\
                      \stmtOpSem{\gamma'_{in+1}}{s'_{n+1}}{stmt^+}{i'_{n+1}, t'_{n+1}, v'_{n+1}}{\gamma'_{on+1}}{s'_{n+2}}{(x'_{n+1}, v'_{n+1})^+}{(i'_{Tn+1}, t'_{Tn+1})^+}.
                  \end{array}
              \end{equation*}

              Because $i'_n$ is not in \importSet{\import{x}, stmt^+}, then
              \Cref{eq:semantics:grfrp:theorems:determinism:sets:import:import} ensures that
              $x \neq i'_n$, so that we have $\gamma'_{in} = \gamma'_n$.
              %
              It implies that:
              \begin{align*}
                  \CurrVal{\tyCtx[x]}{\left[\begin{array}{l}
                                                        \gamma_{i0}[x],   \\
                                                        \dots,            \\
                                                        \gamma_{in-1}[x], \\
                                                        \gamma'_{in}[x],  \\
                                                        \gamma'_{in+1}[x]
                                                    \end{array}\right]} & = \CurrVal{\tyCtx[x]}{\left[\begin{array}{l}
                                                                                                          \text{if} \; x = i_0 \; \text{then} \; v_0 \; \text{else} \; \bot,         \\
                                                                                                          \dots,                                                                     \\
                                                                                                          \text{if} \; x = i_{n-1} \; \text{then} \; v_{n-1} \; \text{else} \; \bot, \\
                                                                                                          \text{if} \; x = i'_n \; \text{then} \; v'_n \; \text{else} \; \bot,       \\
                                                                                                          \text{if} \; x = i'_{n+1} \; \text{then} \; v'_{n+1} \; \text{else} \; \bot
                                                                                                      \end{array}\right]} \\
                                                             & = \CurrVal{\tyCtx[x]}{\left[\begin{array}{l}
                                                                                                       \text{if} \; x = i_0 \; \text{then} \; v_0 \; \text{else} \; \bot,         \\
                                                                                                       \dots,                                                                     \\
                                                                                                       \text{if} \; x = i_{n-1} \; \text{then} \; v_{n-1} \; \text{else} \; \bot, \\
                                                                                                       \bot                                                                       \\
                                                                                                       \text{if} \; x = i_n \; \text{then} \; v_n \; \text{else} \; \bot
                                                                                                   \end{array}\right]}    \\
                                                             & = \CurrVal{\tyCtx[x]}{\left[\begin{array}{l}
                                                                                                       \text{if} \; x = i_0 \; \text{then} \; v_0 \; \text{else} \; \bot, \\
                                                                                                       \dots,                                                             \\
                                                                                                       \text{if} \; x = i_n \; \text{then} \; v_n \; \text{else} \; \bot
                                                                                                   \end{array}\right]}            \\
                                                             & = \CurrVal{\tyCtx[x]}{[\gamma_{i0}[x], \dots, \gamma_{in}[x]]}.
              \end{align*}

              \Cref{eq:semantics:grfrp:theorems:equivalence:deps:import} ensures that
              $\depSet{stmt^+} \subseteq \depSet{\import{x}, stmt^+} \cup \{x\}$. We have as initial
              hypothesis:
              \begin{multline*}
                  \forall y \in \depSet{\import{x}, stmt^+}, \; \left(\Theta[y] = \dType \Rightarrow \gamma'_n[y] = \bot\right) \\
                  \wedge \left(\CurrVal{\tyCtx[y]}{[\gamma_0[y], \dots, \gamma_n[y]]} = \CurrVal{\tyCtx[y]}{[\gamma_0[y], \dots, \gamma_{n-1}[y], \gamma'_n[y], \gamma'_{n+1}[y]]}\right),
              \end{multline*}
              which implies:
              \begin{multline*}
                  \forall y \in \depSet{\import{x}, stmt^+}, \; \left(\Theta[y] = \dType \Rightarrow \gamma'_{in}[y] = \bot\right) \\
                  \wedge \left(\CurrVal{\tyCtx[y]}{[\gamma_{i0}[y], \dots, \gamma_{in}[y]]} = \CurrVal{\tyCtx[y]}{[\gamma_{i0}[y], \dots, \gamma_{in-1}[y], \gamma'_{in}[y], \gamma'_{in+1}[y]]}\right).
              \end{multline*}
              And because $\gamma'_{in} = \gamma'_n$, we shown that:
              \begin{multline*}
                  \gamma'_{in}[x] = \bot \\
                  \wedge \left(\CurrVal{\tyCtx[x]}{[\gamma_{i0}[x], \dots, \gamma_{in}[x]]} = \CurrVal{\tyCtx[x]}{[\gamma_{i0}[x], \dots, \gamma_{in-1}[x], \gamma'_{in}[x], \gamma'_{in+1}[x]]}\right),
              \end{multline*}
              so that in the end we have:
              \begin{multline*}
                  \forall y \in \depSet{stmt^+}, \; \left(\Theta[y] = \dType \Rightarrow \gamma'_{in}[y] = \bot\right) \\
                  \wedge \left(\CurrVal{\tyCtx[y]}{[\gamma_{i0}[y], \dots, \gamma_{in}[y]]} = \CurrVal{\tyCtx[y]}{[\gamma_{i0}[y], \dots, \gamma_{in-1}[y], \gamma'_{in}[y], \gamma'_{in+1}[y]]}\right).
              \end{multline*}

              \Cref{eq:semantics:grfrp:theorems:determinism:sets:import:import} also ensures that
              $i'_n \notin \importSet{stmt^+}$. We can then apply the structural induction
              hypothesis on $stmt^+$, so that we have $[(i'_{Tn}, t'_{Tn})^+] = \emptylist$,
              $[(x'_n, v'_n)^+] = \emptylist$, $[(i'_{Tn+1}, t'_{Tn+1})^+] =
                  [(i_{Tn}, t_{Tn})^+]$, $[(x'_{n+1}, v'_{n+1})^+] = [(x_n, v_n)^+]$,
              $s'_{n+2} = s_{n+1}$, and:
              \begin{multline*}
                  \forall y \in \declSet{stmt^+}, \; \left(\Theta[y] = \dType \Rightarrow \gamma'_{on}[y] = \bot\right) \\
                  \wedge \left(\CurrVal{\tyCtx[y]}{[\gamma_{o0}[y], \dots, \gamma_{on}[y]]} = \CurrVal{\tyCtx[y]}{[\gamma_{o0}[y], \dots, \gamma_{on-1}[y], \gamma'_{on}[y], \gamma'_{on+1}[y]]}\right).
              \end{multline*}

              Because the operational semantics of statements adds new identifiers to the input
              context in order to produce the output context, we have:
              \begin{align*}
                  \gamma_{o0}[x]    & = \gamma_{i0}[x]    &
                  \dots             &                     &
                  \gamma_{on}[x]    & = \gamma_{in}[x]    &
                  \gamma'_{on}[x]   & = \gamma'_{in}[x]   &
                  \gamma'_{on+1}[x] & = \gamma'_{in+1}[x]
              \end{align*}
              \begin{equation*}
                  \CurrVal{\tyCtx[x]}{[\gamma_{o0}[x], \dots, \gamma_{on-1}[x], \gamma'_{on}[x], \gamma'_{on+1}[x]]} = \CurrVal{\tyCtx[x]}{[\gamma_{o0}[x], \dots, \gamma_{on}[x]]}.
              \end{equation*}

              \Cref{eq:semantics:grfrp:theorems:determinism:sets:decl:import} ensures that
              \declSet{\import{x}, stmt^+} equals $\{x\} \cup \declSet{stmt^+}$, and so:
              \begin{multline*}
                  \forall y \in \declSet{\import{x}, stmt^+}, \; \left(\Theta[y] = \dType \Rightarrow \gamma'_{on}[y] = \bot\right) \\
                  \wedge \left(\CurrVal{\tyCtx[y]}{[\gamma_{o0}[y], \dots, \gamma_{on}[y]]} = \CurrVal{\tyCtx[y]}{[\gamma_{o0}[y], \dots, \gamma_{on-1}[y], \gamma'_{on}[y], \gamma'_{on+1}[y]]}\right).
              \end{multline*}

        \item Suppose $stmt^+$ can be decomposed into $\export{x}, stmt^+$.

              The operational rules \nameref{rule:semantics:grfrp:opsem:step:stmt:export:value}
              and \nameref{rule:semantics:grfrp:opsem:step:stmt:export:bottom} ensure that there
              exist output lists $[(x_{i0}, v_{i0})^+]$, \dots, $[(x_{in}, v_{in})^+]$, $[(x'_{in}, v'_{in})^+]$, and
              $[(x'_{in+1}, v'_{in+1})^+]$ such that for all $k \in \{0, \dots, n\}$:
              \begin{equation*}
                  \begin{array}{c}
                      {[(x_k, v_k)^+]} =
                      \left\{
                      \begin{array}{ll}
                          {[(x, v), (x_{ik}, v_{ik})^+]} & \text{if} \; \gamma_k[x] = v    \\
                          {[(x_{ik}, v_{ik})^+]}         & \text{if} \; \gamma_k[x] = \bot
                      \end{array}
                      \right. \\
                      \stmtOpSem{\gamma_k}{s_k}{stmt^+}{i_k, t_k, v_k}{\gamma_{ok}}{s_{k+1}}{(x_{ik}, v_{ik})^+}{(i_{Tk}, t_{Tk})^+}
                  \end{array}
              \end{equation*}
              and:
              \begin{equation*}
                  \begin{array}{c}
                      {[(x'_n, v'_n)^+]} =
                      \left\{
                      \begin{array}{ll}
                          {[(x, v), (x'_{in}, v'_{in})^+]} & \text{if} \; \gamma'_n[x] = v    \\
                          {[(x'_{in}, v'_{in})^+]}         & \text{if} \; \gamma'_n[x] = \bot
                      \end{array}
                      \right.                                                                                                                  \\
                      \stmtOpSem{\gamma'_n}{s_n}{stmt^+}{i'_n, t'_n, v'_n}{\gamma'_{on}}{s'_{n+1}}{(x'_{in}, v'_{in})^+}{(i'_{Tn}, t'_{Tn})^+} \\
                      {[(x'_{n+1}, v'_{n+1})^+]} =
                      \left\{
                      \begin{array}{ll}
                          {[(x, v), (x'_{in+1}, v'_{in+1})^+]} & \text{if} \; \gamma'_{n+1}[x] = v    \\
                          {[(x'_{in+1}, v'_{in+1})^+]}         & \text{if} \; \gamma'_{n+1}[x] = \bot
                      \end{array}
                      \right.                                                                                                                  \\
                      \stmtOpSem{\gamma'_{n+1}}{s'_{n+1}}{stmt^+}{i'_{n+1}, t'_{n+1}, v'_{n+1}}{\gamma'_{on+1}}{s'_{n+2}}{(x'_{in+1}, v'_{in+1})^+}{(i'_{Tn+1}, t'_{Tn+1})^+}.
                  \end{array}
              \end{equation*}

              \Cref{eq:semantics:grfrp:theorems:equivalence:deps:export} ensures that
              \depSet{\export{x}, stmt^+} equals $\{x\} \cup \depSet{stmt^+}$, so that:
              \begin{multline*}
                  \forall y \in \depSet{stmt^+}, \; \left(\Theta[y] = \dType \Rightarrow \gamma'_n[y] = \bot\right) \\
                  \wedge \left(\CurrVal{\tyCtx[y]}{[\gamma_0[y], \dots, \gamma_n[y]]} = \CurrVal{\tyCtx[y]}{[\gamma_0[y], \dots, \gamma_{n-1}[y], \gamma'_n[y], \gamma'_{n+1}[y]]}\right),
              \end{multline*}

              \Cref{eq:semantics:grfrp:theorems:determinism:sets:import:export} ensures that
              $i'_n \notin \importSet{stmt^+}$. We can then apply the structural induction
              hypothesis on $stmt^+$, so that we have $[(i'_{Tn}, t'_{Tn})^+] = \emptylist$,
              $[(x'_{in}, v'_{in})^+] = \emptylist$, $[(i'_{Tn+1}, t'_{Tn+1})^+] =
                  [(i_{Tn}, t_{Tn})^+]$, $[(x'_{in+1}, v'_{in+1})^+] = [(x_{in}, v_{in})^+]$,
              $s'_{n+2} = s_{n+1}$, and:
              \begin{multline*}
                  \forall y \in \declSet{stmt^+}, \; \left(\Theta[y] = \dType \Rightarrow \gamma'_{on}[y] = \bot\right) \\
                  \wedge \left(\CurrVal{\tyCtx[y]}{[\gamma_{o0}[y], \dots, \gamma_{on}[y]]} = \CurrVal{\tyCtx[y]}{[\gamma_{o0}[y], \dots, \gamma_{on-1}[y], \gamma'_{on}[y], \gamma'_{on+1}[y]]}\right).
              \end{multline*}

              \Cref{eq:semantics:grfrp:theorems:determinism:sets:decl:export} ensures that
              \declSet{\export{x}, stmt^+} equals \declSet{stmt^+}, and so:
              \begin{multline*}
                  \forall y \in \declSet{\export{x}, stmt^+}, \; \left(\Theta[y] = \dType \Rightarrow \gamma'_{on}[y] = \bot\right) \\
                  \wedge \left(\CurrVal{\tyCtx[y]}{[\gamma_{o0}[y], \dots, \gamma_{on}[y]]} = \CurrVal{\tyCtx[y]}{[\gamma_{o0}[y], \dots, \gamma_{on-1}[y], \gamma'_{on}[y], \gamma'_{on+1}[y]]}\right).
              \end{multline*}

              The rule \nameref{rule:semantics:grfrp:react_typing:stmt:export} ensures that
              $\Theta[x] = \dType$. Because $x$ is in $\depSet{\export{x}, stmt^+}$, we have
              (initial hypothesis) $\gamma'_n[x] = \bot$. It implies that $[(x_n, v_n)^+]$ equals
              $[(x_{in}, v_{in})^+]$, which is empty.

              The initial hypothesis also ensures that:
              \begin{align*}
                  \CurrVal{\tyCtx[x]}{[\gamma_0[x], \dots, \gamma_n[x]]} & = \CurrVal{\tyCtx[x]}{[\gamma_0[x], \dots, \gamma_{n-1}[x], \gamma'_n[x], \gamma'_{n+1}[x]]} \\
                                                                         & = \CurrVal{\tyCtx[x]}{[\gamma_0[x], \dots, \gamma_{n-1}[x], \bot, \gamma'_{n+1}[x]]}         \\
                                                                         & = \CurrVal{\tyCtx[x]}{[\gamma_0[x], \dots, \gamma_{n-1}[x], \gamma'_{n+1}[x]]}               \\
                                                                         & \Rightarrow \gamma_n[x] = \gamma'_{n+1}[x],
              \end{align*}
              and therefore:

              \begin{align*}
                  {[(x'_{n+1}, v'_{n+1})^+]} & =
                  \left\{
                  \begin{array}{ll}
                      {[(x, v), (x'_{in+1}, v'_{in+1})^+]} & \text{if} \; \gamma'_{n+1}[x] = v    \\
                      {[(x'_{in+1}, v'_{in+1})^+]}         & \text{if} \; \gamma'_{n+1}[x] = \bot
                  \end{array}
                  \right.                                          \\
                                             & =
                  \left\{
                  \begin{array}{ll}
                      {[(x, v), (x_{in}, v_{in})^+]} & \text{if} \; \gamma_n[x] = v    \\
                      {[(x_{in}, v_{in})^+]}         & \text{if} \; \gamma_n[x] = \bot
                  \end{array}
                  \right.                                          \\
                                             & = {[(x_n, v_n)^+]}.
              \end{align*}

        \item Suppose $stmt^+$ can be decomposed into $\decl{x^+}{\sigma}, stmt^+$.

              The rule \nameref{rule:semantics:grfrp:opsem:step:stmt:decl} ensures that
              there exist
              states $s_{\sigma 0}$, \dots, $s_{\sigma n}$, $s'_{\sigma n}$, and $s'_{\sigma n+1}$ corresponding to the flow operation $\sigma$,
              states $s_{i0}$, \dots, $s_{in}$, $s'_{in}$, and $s'_{in+1}$ corresponding to the remaining statements $stmt^+$,
              timer resets $[(i_{\sigma 0}, t_{\sigma 0})^+]$, \dots, $[(i_{\sigma n}, t_{\sigma n})^+]$, $[(i'_{\sigma n}, t'_{\sigma n})^+]$, and $[(i'_{\sigma n+1}, t'_{\sigma n+1})^+]$ from $\sigma$,
              timer resets $[(i_{i0}, t_{i0})^+]$, \dots, $[(i_{in}, t_{in})^+]$, $[(i'_{in}, t'_{in})^+]$, and $[(i'_{in+1}, t'_{in+1})^+]$ from $stmt^+$,
              update environments $\gamma_{i0}$, \dots, $\gamma_{in}$, $\gamma'_{in}$, and $\gamma'_{in+1}$, and
              lists of bottom values $[v_{0\bot}^+]$, \dots, $[v_{n\bot}^+]$, $[v_{n\bot}^{\prime+}]$, and $[v_{n+1\bot}^{\prime+}]$, and
              such that for all $k \in \{0, \dots, n\}$:
              \begin{equation*}
                  \begin{array}{c}
                      s_k = (s_{\sigma k}, s_{ik})                                                                                                                \quad
                      \gamma_{ik} = \gamma_k \ctxconcat [(v_{k\bot}/x)^+]                                                                                         \quad
                      { [(i_{Tk}, t_{Tk})^+] = [(i_{\sigma k}, t_{\sigma k})^+, (i_{ik}, t_{ik})^+] }                                  \\
                      \sigmaOpSem{\gamma_k}{s_{\sigma k}}{\sigma}{t_k}{s_{\sigma  k+1}}{v_{k\bot}^+}{(i_{\sigma k}, t_{\sigma k})^+}   \\
                      \stmtOpSem{\gamma_{ik}}{s_{ik}}{stmt^+}{i_k, t_k, v_k}{\gamma_{ok}}{s_{i k+1}}{(x_k, v_k)^+}{(i_{ik}, t_{ik})^+} \\
                  \end{array}
              \end{equation*}
              and:
              \begin{equation*}
                  \begin{array}{c}
                      s'_n = (s'_{\sigma n}, s'_{in})                                                                                                             \quad
                      \gamma'_{in} = \gamma'_n \ctxconcat [(v'_{n\bot}/x)^+]                                                                                      \quad
                      { [(i'_{Tn}, t'_{Tn})^+] = [(i'_{\sigma n}, t'_{\sigma n})^+, (i'_{in}, t'_{in})^+] }                                                       \\
                      \sigmaOpSem{\gamma'_n}{s'_{\sigma n}}{\sigma}{t'_n}{s'_{\sigma n+1}}{v_{n\bot}^{\prime+}}{(i'_{\sigma n}, t'_{\sigma n})^+}                 \\
                      \stmtOpSem{\gamma'_{in}}{s'_{in}}{stmt^+}{i'_n, t'_n, v'_n}{\gamma'_{on}}{s'_{in+1}}{(x'_n, v'_n)^+}{(i'_{in}, t'_{in})^+}                  \\
                      s'_{n+1} = (s'_{\sigma n+1}, s'_{in+1})                                                                                                     \quad
                      \gamma'_{in+1} = \gamma'_{n+1} \ctxconcat [(v'_{n+1\bot}/x)^+]                                                                              \\
                      { [(i'_{Tn+1}, t'_{Tn+1})^+] = [(i'_{\sigma n+1}, t'_{\sigma n+1})^+, (i'_{in+1}, t'_{in+1})^+] }                                           \\
                      \sigmaOpSem{\gamma'_{n+1}}{s'_{\sigma n+1}}{\sigma}{t'_{n+1}}{s'_{\sigma n+2}}{v_{n+1\bot}^{\prime+}}{(i'_{\sigma n+1}, t'_{\sigma n+1})^+} \\
                      \stmtOpSem{\gamma'_{in+1}}{s'_{in+1}}{stmt^+}{i'_{n+1}, t'_{n+1}, v'_{n+1}}{\gamma'_{on+1}}{s'_{in+2}}{(x'_{n+1}, v'_{n+1})^+}{(i'_{in+1}, t'_{in+1})^+}.
                  \end{array}
              \end{equation*}

              The typing rules \nameref{rule:grust:typing:grfrp:flow_stmt:let} and
              \nameref{rule:grust:typing:grfrp:flow_stmt:def} ensures that we have:
              \begin{equation*}
                  \flowTyped{\sigma}{(\tyCtx[x])^+}.
              \end{equation*}

              \Cref{eq:semantics:grfrp:theorems:equivalence:deps:decl} ensures that
              \depSet{\sigma} is included in \depSet{\decl{x^+}{\sigma}, stmt^+}, so that:
              \begin{multline*}
                  \forall y \in \depSet{\sigma}, \; \left(\Theta[y] = \dType \Rightarrow \gamma'_n[y] = \bot\right) \\
                  \wedge \left(\CurrVal{\tyCtx[y]}{[\gamma_0[y], \dots, \gamma_n[y]]} = \CurrVal{\tyCtx[y]}{[\gamma_0[y], \dots, \gamma_{n-1}[y], \gamma'_n[y], \gamma'_{n+1}[y]]}\right).
              \end{multline*}
              We can then apply \Cref{lem:semantics:grfrp:theorems:oversampling:op}, which ensures:
              \begin{equation*}
                  \begin{array}{c}
                      [(i'_{\sigma n}, t'_{\sigma n})^+] = \emptylist \quad [(i'_{\sigma n+1}, t'_{\sigma n+1})^+] = [(i_{\sigma n}, t_{\sigma n})^+] \\
                      s'_{\sigma n+2} = s_{\sigma n+1} \quad \left(\Theta[x] = \dType \Rightarrow v'_{n\bot} = \bot\right)^+                          \\
                      \left(\CurrVal{\tyCtx[x]}{[v_{0\bot}, \dots, v_{n\bot}]} = \CurrVal{\tyCtx[x]}{[v_{0\bot}, \dots, v_{n-1\bot}, v'_{n\bot}, v'_{n+1\bot}]}\right)^+,
                  \end{array}
              \end{equation*}
              thus we have:
              \begin{multline}\label{eq:semantics:grfrp:theorems:oversampling:stmt:decl:defs}
                  \forall y \in \{x^+\}, \; \Theta[y] = \dType \Rightarrow \gamma'_{in}[y] = \bot \\
                  \wedge \left(\CurrVal{\tyCtx[y]}{[\gamma_{i0}[y], \dots, \gamma_{in}[y]]} = \CurrVal{\tyCtx[y]}{[\gamma_{i0}[y], \dots, \gamma_{in-1}[y], \gamma'_{in}[y], \gamma'_{in+1}[y]]}\right),
              \end{multline}

              \Cref{eq:semantics:grfrp:theorems:equivalence:deps:decl} also ensures that
              $\depSet{stmt^+} \subseteq \depSet{\decl{x^+}{\sigma}, stmt^+} \cup \{x^+\}$. We have
              as initial hypothesis:
              \begin{multline*}
                  \forall y \in \depSet{\decl{x^+}{\sigma}, stmt^+}, \; \left(\Theta[y] = \dType \Rightarrow \gamma'_n[y] = \bot\right) \\
                  \wedge \left(\CurrVal{\tyCtx[y]}{[\gamma_0[y], \dots, \gamma_n[y]]} = \CurrVal{\tyCtx[y]}{[\gamma_0[y], \dots, \gamma_{n-1}[y], \gamma'_n[y], \gamma'_{n+1}[y]]}\right),
              \end{multline*}
              which implies:
              \begin{multline*}
                  \forall y \in \depSet{\decl{x^+}{\sigma}, stmt^+}, \; \left(\Theta[y] = \dType \Rightarrow \gamma'_{in}[y] = \bot\right) \\
                  \wedge \left(\CurrVal{\tyCtx[y]}{[\gamma_{i0}[y], \dots, \gamma_{in}[y]]} = \CurrVal{\tyCtx[y]}{[\gamma_{i0}[y], \dots, \gamma_{in-1}[y], \gamma'_{in}[y], \gamma'_{in+1}[y]]}\right).
              \end{multline*}
              so that, with \Cref{eq:semantics:grfrp:theorems:oversampling:stmt:decl:defs}, we have:
              \begin{multline*}
                  \forall y \in \depSet{stmt^+}, \; \left(\Theta[y] = \dType \Rightarrow \gamma'_{in}[y] = \bot\right) \\
                  \wedge \left(\CurrVal{\tyCtx[y]}{[\gamma_{i0}[y], \dots, \gamma_{in}[y]]} = \CurrVal{\tyCtx[y]}{[\gamma_{i0}[y], \dots, \gamma_{in-1}[y], \gamma'_{in}[y], \gamma'_{in+1}[y]]}\right).
              \end{multline*}

              \Cref{eq:semantics:grfrp:theorems:determinism:sets:import:decl} also ensures that
              $i'_n \notin \importSet{stmt^+}$. We can then apply the structural induction
              hypothesis on $stmt^+$, so that we have $[(i'_{in}, t'_{in})^+] = \emptylist$,
              $[(x'_n, v'_n)^+] = \emptylist$, $[(i'_{in+1}, t'_{in+1})^+] =
                  [(i_{in}, t_{in})^+]$, $[(x'_{n+1}, v'_{n+1})^+] = [(x_n, v_n)^+]$,
              $s'_{in+2} = s_{in+1}$, and:
              \begin{multline*}
                  \forall y \in \declSet{stmt^+}, \; \left(\Theta[y] = \dType \Rightarrow \gamma'_{on}[y] = \bot\right) \\
                  \wedge \left(\CurrVal{\tyCtx[y]}{[\gamma_{o0}[y], \dots, \gamma_{on}[y]]} = \CurrVal{\tyCtx[y]}{[\gamma_{o0}[y], \dots, \gamma_{on-1}[y], \gamma'_{on}[y], \gamma'_{on+1}[y]]}\right).
              \end{multline*}

              Because the operational semantics of statements adds new identifiers to the input
              context in order to produce the output context,
              for all defined identifier $x$ we have:
              \begin{align*}
                  \gamma_{o0}[x]    & = \gamma_{i0}[x]    &
                  \dots             &                     &
                  \gamma_{on}[x]    & = \gamma_{in}[x]    &
                  \gamma'_{on}[x]   & = \gamma'_{in}[x]   &
                  \gamma'_{on+1}[x] & = \gamma'_{in+1}[x]
              \end{align*}
              \begin{equation*}
                  \CurrVal{\tyCtx[x]}{[\gamma_{o0}[x], \dots, \gamma_{on-1}[x], \gamma'_{on}[x], \gamma'_{on+1}[x]]} = \CurrVal{\tyCtx[x]}{[\gamma_{o0}[x], \dots, \gamma_{on}[x]]}.
              \end{equation*}

              \Cref{eq:semantics:grfrp:theorems:determinism:sets:decl:decl} states that
              \declSet{\decl{x^+}{\sigma}, stmt^+} equals $\declSet{stmt^+} \cup {x^+}$, and so:
              \begin{multline*}
                  \forall y \in \declSet{\decl{x^+}{\sigma}, stmt^+}, \; \left(\Theta[y] = \dType \Rightarrow \gamma'_{on}[y] = \bot\right) \\
                  \wedge \left(\CurrVal{\tyCtx[y]}{[\gamma_{o0}[y], \dots, \gamma_{on}[y]]} = \CurrVal{\tyCtx[y]}{[\gamma_{o0}[y], \dots, \gamma_{on-1}[y], \gamma'_{on}[y], \gamma'_{on+1}[y]]}\right).
              \end{multline*}
    \end{itemize}
\end{skippableproof}

\begin{theorem}[Oversampling]
    \label{thm:semantics:grfrp:theorems:oversampling}
    Let \service{F}{stmt^+} be a service in a program \G with a typing context \tyCtx, and let
    $\Theta$ be a reactive typing context such that:
    \begin{equation*}
        \itemTyped{\service{F}{stmt^+}} \qquad \servTypedReact{\service{F}{stmt^+}}.
    \end{equation*}
    For all $n \ge 1$;
    all states $s_0$, \dots, $s_{n+1}$, $s'_{n+1}$, and $s'_{n+2}$ corresponding to $F$;
    all timer queues $\tau_0$, \dots, $\tau_n$, $\tau'_{n+1}$, and $\tau'_{n+2}$;
    all messages $\mess_0$, \dots, $\mess_n$, $\mess'_n$, and $\mess'_{n+1}$; and
    all timestamps $t_0$, \dots, $t_n$, $t'_n$, and $t'_{n+1}$;
    all lists of outputs $[(x_0, v_0)^+]$, \dots, $[(x_n, v_n)^+]$, $[(x'_n, v'_n)^+]$, and $[(x'_{n+1}, v'_{n+1})^+]$
    such that:
    \begin{equation*}
        \begin{array}{c}
            t'_n \le t'_{n+1} \qquad (\mess'_{n+1}, t'_{n+1}) = (\mess_n, t_n) \qquad s_0 = \initstate{F} \qquad \tau_0 = \inittimer{F} \\
            \forall k \in \{0, \dots, n\}, \quad \serviceOpSem{s_k}{\tau_k}{F}{\mess_k}{t_k}{s_1}{\tau_{k+1}}{(x_k, v_k)^+}             \\
            \serviceOpSem{s_n}{\tau_n}{F}{\mess'_n}{t_n}{s'_{n+1}}{\tau'_{n+1}}{(x'_n, v'_n)^+}                                         \\
            \serviceOpSem{s_{n+1}}{\tau_{n+1}}{F}{\mess'_{n+1}}{t_{n+1}}{s'_{n+2}}{\tau'_{n+2}}{(x'_{n+1}, v'_{n+1})^+},
        \end{array}
    \end{equation*}
    if there exist an identifier $i \notin \importSet{F}$ and a value $v_i$ such that
    $\mess'_n = \inputMess{i}{v_i}$, then we have: % \todo{unknown timer message}
    \begin{align*}
        \tau'_{n+1} & = \tau_n     &          &           & [(x'_n, v'_n)^+]         & = \emptylist      \\
        \tau'_{n+2} & = \tau_{n+1} & s'_{n+2} & = s_{n+1} & [(x'_{n+1}, v'_{n+1})^+] & = [(x_n, v_n)^+].
    \end{align*}

    This lemma formalizes the behavior of a service, $F$, when it receives an input that it does not
    import. It states that if, after $n$ simulation steps, $F$ receives such an input at step $n+1$,
    this intermediate step has no lasting effect. A computation performed from this new state
    produces the same result as one performed from the state at step $n$. This implies that the
    internal state of the service does not change ($s'_{n+2} = s_{n+1}$), the timer queues are
    unaffected ($\tau'_{n+1} = \tau_n$ and $\tau'_{n+2} = \tau_{n+1}$), and the exported values are
    preserved ($[(x'_n, v'_n)^+] = \emptylist$ and $[(x'_{n+1}, v'_{n+1})^+] = [(x_n, v_n)^+]$).
\end{theorem}
\begin{skippableproof}
    Let \service{F}{stmt^+} be a service in a program \G with a typing context \tyCtx, let
    $\Theta$ be a reactive typing context such that:
    \begin{equation*}
        \itemTyped{\service{F}{stmt^+}} \qquad \servTypedReact{\service{F}{stmt^+}}.
    \end{equation*}
    Let $n \ge 1$, let $s_0$, \dots, $s_{n+1}$, $s'_{n+1}$, and $s'_{n+2}$ be states corresponding
    to $F$, let $\tau_0$, \dots, $\tau_n$, $\tau'_{n+1}$, and $\tau'_{n+2}$ be timer
    queues, let $\mess_0$, \dots, $\mess_n$, $\mess'_n$, and $\mess'_{n+1}$ be messages, and let
    $[(x_0, v_0)^+]$, \dots, $[(x_n, v_n)^+]$, $[(x'_n, v'_n)^+]$, and $[(x'_{n+1}, v'_{n+1})^+]$ be
    lists of outputs such that:
    \begin{equation*}
        \begin{array}{c}
            t'_n \le t'_{n+1} = t_n \qquad s_0 = \initstate{F} \qquad \tau_0 = \inittimer{F}                                \\
            \forall k \in \{0, \dots, n\}, \quad \serviceOpSem{s_k}{\tau_k}{F}{\mess_k}{t_k}{s_1}{\tau_{k+1}}{(x_k, v_k)^+} \\
            \serviceOpSem{s_n}{\tau_n}{F}{\mess'_n}{t_n}{s'_{n+1}}{\tau'_{n+1}}{(x'_n, v'_n)^+}                             \\
            \serviceOpSem{s_{n+1}}{\tau_{n+1}}{F}{\mess'_{n+1}}{t_{n+1}}{s'_{n+2}}{\tau'_{n+2}}{(x'_{n+1}, v'_{n+1})^+}.
        \end{array}
    \end{equation*}

    Suppose there exist an identifier $i$ not in \importSet{F} and a value $v_i$ such that
    $\mess'_n$ equals \inputMess{i}{v_i}.

    The operational semantics rules \nameref{rule:semantics:grfrp:opsem:step:service:input} and
    \nameref{rule:semantics:grfrp:opsem:step:service:timer} ensure there exist
    input update contexts $\gamma_{i0}$, \dots, $\gamma_{in}$, $\gamma'_{in}$, and $\gamma'_{in+1}$;
    output update contexts $\gamma_{o0}$, \dots, $\gamma_{on}$, $\gamma'_{on}$, and $\gamma'_{on+1}$;
    triggering inputs $(i_0, t_0, v_0)$, \dots, $(i_n, t_n, v_n)$, $(i'_n, t'_n, v'_n)$, and $(i'_{n+1}, t'_{n+1}, v'_{n+1})$;
    lists of outputs $[(x_0, v_0)^+]$, \dots, $[(x_n, v_n)^+]$, $[(x'_n, v'_n)^+]$, and $[(x'_{n+1}, v'_{n+1})^+]$; and
    lists of timer resets $[(i_{T0}, t_{T0})^+]$, \dots, $[(i_{Tn}, t_{Tn})^+]$, $[(i'_{Tn}, t'_{Tn})^+]$, and $[(i'_{Tn+1}, t'_{Tn+1})^+]$
    such that for all $k \in \{0, \dots, n\}$:
    \begin{equation*}
        \begin{array}{c}
            (\gamma_{ik}, v_k, \tau_{k+1}) =
            \left\{
            \begin{array}{lll}
                ([()/i_k],    & (),  & \reset{\tau}{[(i_{Tk}, t_{Tk})^+]})                                                           \\
                              &      & \hspace{-1cm} \text{if} \; \mess_k = \timerMess{i_k} \wedge \tau_k = (i_k, t_k) \diamond \tau \\
                (\emptygamma, & v_i, & \reset{\tau_k}{[(i_{Tk}, t_{Tk})^+]})                                                         \\
                              &      & \hspace{-1cm} \text{if} \; \mess_k = \inputMess{i_k}{v_i}
            \end{array}
            \right. \\
            \stmtOpSem{\gamma_{ik}}{s_k}{stmt^+}{i_k, t_k, v_k}{\gamma_{ok}}{s_{k+1}}{(x_k, v_k)^+}{(i_{Tk}, t_{Tk})^+}
        \end{array}
    \end{equation*}
    and:
    \begin{equation*}
        \begin{array}{c}
            (\gamma'_{in}, v'_n, \tau'_{n+1}) =
            \left\{
            \begin{array}{lll}
                ([()/i'_n],   & (),  & \reset{\tau}{[(i'_{Tn}, t'_{Tn})^+]})                                                             \\
                              &      & \hspace{-1cm} \text{if} \; \mess'_n = \timerMess{i'_n} \wedge \tau_n = (i'_n, t'_n) \diamond \tau \\
                (\emptygamma, & v_i, & \reset{\tau_n}{[(i'_{Tn}, t'_{Tn})^+]})                                                           \\
                              &      & \hspace{-1cm} \text{if} \; \mess'_n = \inputMess{i'_n}{v_i}
            \end{array}
            \right.                                                                                                               \\
            \stmtOpSem{\gamma'_{in}}{s_n}{stmt^+}{i'_n, t'_n, v'_n}{\gamma'_{on}}{s'_{n+1}}{(x'_n, v'_n)^+}{(i'_{Tn}, t'_{Tn})^+} \\
            (\gamma'_{in+1}, v'_{n+1}, \tau'_{n+2}) =
            \left\{
            \begin{array}{lll}
                ([()/i'_{n+1}], & (),  & \reset{\tau}{[(i'_{Tn+1}, t'_{Tn+1})^+]})                                                                              \\
                                &      & \hspace{-1cm} \text{if} \; \mess'_{n+1} = \timerMess{i'_{n+1}} \wedge \tau'_{n+1} = (i'_{n+1}, t'_{n+1}) \diamond \tau \\
                (\emptygamma,   & v_i, & \reset{\tau'_{n+1}}{[(i'_{Tn+1}, t'_{Tn+1})^+]})                                                                       \\
                                &      & \hspace{-1cm} \text{if} \; \mess'_{n+1} = \inputMess{i'_{n+1}}{v_i}
            \end{array}
            \right.                                                                                                               \\
            \stmtOpSem{\gamma'_{in+1}}{s'_{n+1}}{stmt^+}{i'_{n+1}, t'_{n+1}, v'_{n+1}}{\gamma'_{on+1}}{s'_{n+2}}{(x'_{n+1}, v'_{n+1})^+}{(i'_{Tn+1}, t'_{Tn+1})^+}.
        \end{array}
    \end{equation*}

    We supposed that $\mess'_n = \inputMess{i}{v_i}$, with $i \notin \importSet{F}$, which implies:
    \begin{align*}
        \gamma'_{in} & = \emptygamma & v'_n & = v_i & \tau'_{n+1} & = \reset{\tau_n}{[(i'_{Tn}, t'_{Tn})^+]}.
    \end{align*}

    \Cref{eq:semantics:grfrp:theorems:determinism:sets:import:service} states that \importSet{F}
    equals \importSet{stmt^+}.

    Let $x \in \depSet{stmt^+}$, then we have:
    \begin{align*}
        \gamma'_{in}[x]   & = \emptygamma[x] = \bot \\
        \gamma'_{in+1}[x] & =
        \left\{
        \begin{array}{ll}
            [()/i'_{n+1}][x] & \text{if} \; \mess'_{n+1} = \timerMess{i'_{n+1}} \wedge \tau'_{n+1} = (i'_{n+1}, t'_{n+1}) \diamond \tau \\
            \emptygamma[x]   & \text{if} \; \mess'_{n+1} = \inputMess{i'_{n+1}}{v_i}
        \end{array}
        \right.                                     \\
                          & =
        \left\{
        \begin{array}{ll}
            \bot & \text{if} \; \mess'_{n+1} = \timerMess{i'_{n+1}} \wedge \tau'_{n+1} = (i'_{n+1}, t'_{n+1}) \diamond \tau \\
            \bot & \text{if} \; \mess'_{n+1} = \inputMess{i'_{n+1}}{v_i}
        \end{array}
        \right.                                     \\
                          & = \bot                  \\
        \gamma_{in}[x]    & =
        \left\{
        \begin{array}{ll}
            [()/i_{n}][x]  & \text{if} \; \mess_{n} = \timerMess{i_{n}} \wedge \tau_{n} = (i_{n}, t_{n}) \diamond \tau \\
            \emptygamma[x] & \text{if} \; \mess_{n} = \inputMess{i_{n}}{v_i}
        \end{array}
        \right.                                     \\
                          & = \bot
    \end{align*}
    \begin{align*}
        \CurrVal{\tyCtx[x]}{[\gamma_{i0}[x], \dots, \gamma_{in-1}[x], \gamma'_{in}[x], \gamma'_{in+1}[x]]} & = \CurrVal{\tyCtx[x]}{[\gamma_{i0}[x], \dots, \gamma_{in-1}[x], \bot, \bot]}      \\
                                                                                                           & = \CurrVal{\tyCtx[x]}{[\gamma_{i0}[x], \dots, \gamma_{in-1}[x], \gamma_{in}[x]]},
    \end{align*}
    which implies:
    \begin{multline*}
        \forall x \in \depSet{stmt^+}, \; \left(\Theta[x] = \dType \Rightarrow \gamma'_{in}[x] = \bot\right) \\
        \wedge \left(\CurrVal{\tyCtx[x]}{[\gamma_{i0}[x], \dots, \gamma_{in}[x]]} = \CurrVal{\tyCtx[x]}{[\gamma_{i0}[x], \dots, \gamma_{in-1}[x], \gamma'_{in}[x], \gamma'_{in+1}[x]]}\right).
    \end{multline*}

    The application of \Cref{lem:semantics:grfrp:theorems:oversampling:stmts} ensures that:
    \begin{multline*}
        \forall y \in \declSet{stmt^+}, \; \left(\Theta[y] = \dType \Rightarrow \gamma'_{on}[y] = \bot\right) \\
        \wedge \left(\CurrVal{\tyCtx[y]}{[\gamma_{o0}[y], \dots, \gamma_{on}[y]]} = \CurrVal{\tyCtx[y]}{[\gamma_{o0}[y], \dots, \gamma_{on-1}[y], \gamma'_{on}[y], \gamma'_{on+1}[y]]}\right)
    \end{multline*}
    and:
    \begin{align*}
        [(i'_{Tn}, t'_{Tn})^+]     & = \emptylist           &          &           & [(x'_n, v'_n)^+]         & = \emptylist      \\
        [(i'_{Tn+1}, t'_{Tn+1})^+] & = [(i_{Tn}, t_{Tn})^+] & s'_{n+2} & = s_{n+1} & [(x'_{n+1}, v'_{n+1})^+] & = [(x_n, v_n)^+].
    \end{align*}

    Finally, because we supposed $\mess'_{n+1} = \mess_n$, we have:
    \begin{align*}
        \tau'_{n+1} & = \reset{\tau_n}{[(i'_{Tn}, t'_{Tn})^+]} = \reset{\tau_n}{\emptylist} = \tau_n \\
        \tau'_{n+2} & =
        \left\{
        \begin{array}{lll}
            \reset{\tau}{[(i'_{Tn+1}, t'_{Tn+1})^+]}        & \text{if} \; \mess'_{n+1} = \timerMess{i'_{n+1}}              \\
                                                            & \quad \wedge \tau'_{n+1} = (i'_{n+1}, t'_{n+1}) \diamond \tau \\
            \reset{\tau'_{n+1}}{[(i'_{Tn+1}, t'_{Tn+1})^+]} & \text{if} \; \mess'_{n+1} = \inputMess{i'_{n+1}}{v_i}
        \end{array}
        \right.                                                                                      \\
                    & =
        \left\{
        \begin{array}{lll}
            \reset{\tau}{[(i_{Tn}, t_{Tn})^+]}   & \text{if} \; \mess_n = \timerMess{i_n} \wedge \tau_n = (i_n, t_n) \diamond \tau \\
            \reset{\tau_n}{[(i_{Tn}, t_{Tn})^+]} & \text{if} \; \mess_n = \inputMess{i_n}{v_i}
        \end{array}
        \right.                                                                                      \\
                    & = \tau_{n+1}.
    \end{align*}
\end{skippableproof}

\paragraph{Summary}
In conclusion, this section established the oversampling property, a cornerstone of the \GRust
execution model. The main theorem, \Cref{thm:semantics:grfrp:theorems:oversampling}, formally proves
that a service's behavior is driven solely by changes to its imported flows.
%
By demonstrating that an evaluation triggered by an unimported flow---an oversampling step---%
produces no outputs and leaves the service state and timers unchanged, we confirm that such
computations are inert.
%
The proof was built upon a key lemma for reactive flow operations, which required introducing a
history-aware \smf{curr} function to formally compare flow values across computation steps.
%
This oversampling guarantee is crucial: it provides the formal justification for adopting an
efficient, event-driven execution model for services without compromising the semantic integrity of
the underlying reactive components.
